% Created 2026-08-02 dom 19:30
% Intended LaTeX compiler: pdflatex
\documentclass[14pt, a4paper]{article}
\usepackage[utf8]{inputenc}
\usepackage[T1]{fontenc}
\usepackage{graphicx}
\usepackage{longtable}
\usepackage{wrapfig}
\usepackage{rotating}
\usepackage[normalem]{ulem}
\usepackage{amsmath}
\usepackage{amssymb}
\usepackage{capt-of}
\usepackage{hyperref}
\usepackage[utf8]{inputenc}
\usepackage[T1]{fontenc}
\usepackage[spanish, es-tabla]{babel}
\usepackage{libertine}
\usepackage{libertinust1math}
\usepackage{inconsolata}
\usepackage{textalpha}
\usepackage{textcomp}
\usepackage{geometry}
\geometry{a4paper, left=3.5cm, right=3cm, top=3cm, bottom=3cm}
\usepackage{fancyhdr}
\pagestyle{fancy}
\fancyhf{}
\fancyhead[LE,RO]{\large\thepage}
\fancyhead[RE]{\large\nouppercase{\leftmark}}
\fancyhead[LO]{\large\nouppercase{\rightmark}}
\renewcommand{\headrulewidth}{0.4pt}
\usepackage{titlesec}
\titleformat{\section}{\normalfont\Huge\bfseries}{\thesection}{1em}{}
\titleformat{\subsection}{\normalfont\LARGE\bfseries}{\thesubsection}{1em}{}
\titleformat{\subsubsection}{\normalfont\Large\bfseries}{\thesubsubsection}{1em}{}
\usepackage{setspace}
\onehalfspacing
\usepackage{indentfirst}
\setlength{\parindent}{1.5em}
\setlength{\parskip}{0.6em plus 0.2em minus 0.2em}
\usepackage{microtype}
\usepackage[colorlinks=true,linkcolor=black,citecolor=black,urlcolor=black]{hyperref}
\hyphenation{elo-gio}
\usepackage[spanish]{babel}
\selectlanguage{spanish}
\author{Rigel Valencia Osorio, hijo de Beatriz y Guillermo}
\date{Julio de 2026}
\title{Hacia una Ontología Materialista de la Música desde la Experiencia Vital\\\medskip
\large Filosofía, sonido y soberanía tecnológica — Postulación como Sabedor ante el Grupo Tántalo}
\hypersetup{
 pdfauthor={Rigel Valencia Osorio, hijo de Beatriz y Guillermo},
 pdftitle={Hacia una Ontología Materialista de la Música desde la Experiencia Vital},
 pdfkeywords={},
 pdfsubject={},
 pdfcreator={Emacs 30.1 (Org mode 9.7.29)}, 
 pdflang={Spanish}}
\begin{document}

\maketitle
\tableofcontents

\section{Resumen Ejecutivo}
\label{sec:org8c2b344}
El presente documento constituye una postulación como Sabedor ante el
Grupo de Investigación Tántalo del Departamento de Filosofía de la
Universidad de Caldas, figura contemplada en los reglamentos internos
de la institución para el reconocimiento de saberes adquiridos fuera
del circuito académico formal, en consonancia con la figura del
Experto definida por la Ley 30 de 1992 y la Sentencia C-109 de 1994
de la Corte Constitucional. Se inscribe en la intersección entre la
filosofía de la música, la ontología materialista y la praxis
tecnológica mediante software libre, proponiendo una visión de la
música como laboratorio filosófico donde se verifican tesis sobre la
discontinuidad de la materia, el entrelazamiento de géneros
(\textbf{symploké}) y la crítica al idealismo subjetivista.

Se expone una trayectoria formativa que constituye una dialéctica
constante entre la academia formal y el conocimiento empírico. A
partir de los conceptos de "paradigma" de Thomas Kuhn, "cierre
categorial" de Gustavo Bueno y "elogio de la dificultad" de
Estanislao Zuleta, se argumenta que la interrupción de los estudios
formales no implica una fractura epistemológica, sino una ampliación
del campo de acción musical y filosófico.

Este ensayo dialoga directamente con las líneas de investigación en
Estética y Filosofía del Arte del Grupo Tántalo. La praxis con
software libre no es aquí un añadido técnico, sino una metodología de
investigación filosófica. No trabajo con una herramienta: trabajo con
una ecología de herramientas que se entrelazan —SuperCollider, Csound,
FoxDot, Hydrogen, Ardour, LSP, AMS, JackControl, Emacs/Lisp—,
modificando y recombinando código libre, explorando la tesitura y las
capacidades tímbricas de cada programa, enrutando audio entre
aplicaciones mientras la guitarra eléctrica suena sobre el código en
tiempo real. La soberanía tecnológica no es una consecuencia
automática del materialismo filosófico, sino una elección consciente
dentro de un campo de fuerzas donde el software privativo y el
software libre compiten. Yo elijo el software libre. No porque sea
inevitable, sino porque es coherente con una ontología que afirma que
el conocimiento no tiene dueño.

No pido un título: pido que se reconozca, bajo la figura del Sabedor
que contemplan los reglamentos de la Universidad de Caldas, un saber
que ya existe y que puede contribuir a la vida filosófica del
Departamento.


\begin{quote}
"La dificultad del camino recorrido no es un obstáculo, sino el método
mismo del conocimiento auténtico."
— Estanislao Zuleta
\end{quote}
\section{Defensa Metodológica: La Investigación-Acción Participativa como Fundamento}
\label{sec:org7603796}
Mi presencia en los territorios del Ariari, en el Distrito de
Aguablanca, en los ensambles de jazz de la Universidad INCA y en las
comunidades de software libre no ha sido la de un observador externo
que extrae datos para luego analizarlos desde la comodidad de un
cubículo académico. Ha sido la de un sujeto que se sumerge en la
realidad para transformarla y ser transformado por ella. Esta
posición no es un capricho biográfico; es una decisión metodológica
que se inscribe en la tradición de la Investigación-Acción
Participativa (IAP), consolidada por el sociólogo colombiano Orlando
Fals Borda.
\subsection{La IAP como metodología del Sur Global}
\label{sec:org26ed332}
La IAP, nacida en el Sur Global como respuesta crítica a las
metodologías extractivistas del Norte, sostiene que el conocimiento
no se produce en la distancia aséptica entre un sujeto que investiga
y un objeto que es investigado, sino en la relación dialéctica entre
ambos. El investigador no es un espectador; es un participante. La
comunidad no es un yacimiento de datos; es una interlocutora. Y el
saber que emerge de esa relación no es propiedad del académico, sino
patrimonio del colectivo que lo ha producido.

Desde el materialismo filosófico, esta posición es perfectamente
coherente. La IAP es, en el fondo, una metodología de la \textbf{symploké}
(συμπλοκή): entrelaza al investigador con la comunidad, al saber
académico con el saber popular, a la teoría con la práctica, sin
confundirlos. No es que el investigador "baje" al territorio a
iluminar a los ignorantes; es que el territorio interpela al
investigador y lo obliga a replantear sus categorías. Como aprendí
con los músicos naturales: mi tarea no era enseñarles, sino encontrar
estrategias para que ellos tuvieran herramientas sin perder su
esencia. Ellos me enseñaron tanto como yo pude ofrecerles.
\subsection{Fals Borda: ciencia y liberación}
\label{sec:org74887d9}
Orlando Fals Borda (1925-2008), sociólogo colombiano, es una de las
figuras más influyentes del pensamiento latinoamericano. Fundador de
la IAP junto con Paulo Freire y otros intelectuales del Sur, Fals
Borda insistía en que la investigación no puede ser neutral: o está
al servicio de la dominación, o está al servicio de la liberación. Su
obra, que combinaba el rigor sociológico con el compromiso político,
demostró que es posible hacer ciencia sin renunciar a la militancia,
y que el conocimiento producido por las comunidades campesinas,
indígenas y afrodescendientes tiene tanto valor como el producido en
las universidades.

Fals Borda decía que "la ciencia no es un fetiche, sino una
herramienta de liberación". Eso es lo que este ensayo pretende: usar
las herramientas de la filosofía materialista, de la musicología, de
la física y de la computación, no para adornar un currículum, sino
para liberar un saber que ha sido producido en los márgenes de la
academia y que merece ser reconocido en su centro.
\subsection{Blindaje frente a dos críticas}
\label{sec:org720daf3}
Esta metodología me blinda frente a dos críticas que la academia
podría hacer a mi ensayo. La primera es la del "extractivismo
epistemológico": ¿con qué derecho viene este músico sin título a
hablarnos de los saberes de las periferias? La respuesta es: con el
derecho que da haber estado allí, no como turista, sino como
compañero. La segunda es la del "subjetivismo": ¿cómo podemos confiar
en un conocimiento que se basa en la experiencia personal? La
respuesta es: la IAP no es subjetivismo; es intersubjetividad. El
conocimiento que emerge de la IAP no es la opinión de un individuo;
es el producto de una relación dialéctica entre el investigador y la
comunidad, validado por la práctica y orientado a la transformación.
\subsection{La IAP y la figura del Sabedor}
\label{sec:orgb9b2e66}
La figura del "Sabedor" —contemplada en los reglamentos internos de la
Universidad de Caldas, no en la Ley 30 de 1992, pero en consonancia
con la figura del "Experto" definida por la Corte Constitucional en la
Sentencia C-109 de 1994— encuentra en la IAP su fundamento
metodológico. Si el "Sabedor" es aquella persona que, sin título
profesional, acredita una preparación especial, empírica o técnica,
entonces la IAP es el método que permite acreditar esa preparación: no
mediante papeles, sino mediante la práctica transformadora. Mi
trayectoria no es una acumulación de experiencias inconexas; es una
investigación-acción que ha producido conocimientos musicales,
pedagógicos y tecnológicos en contextos donde la academia no llega. Lo
que solicito no es un privilegio; es el reconocimiento de una
metodología que la propia academia latinoamericana ha validado.
\section{Introducción: El Músico como Sujeto Ontológico, Político y Digital}
\label{sec:org5e8fb7d}
El músico no es un producto de la industria y el comercio; tampoco
niega su valía. Esta es la primera afirmación que sostiene este
ensayo, y quizás la más incómoda para una sociedad que ha reducido el
arte a entretenimiento o a mercancía. El músico es, ante todo, un ser
político activo: consciente de su posición en la estructura social y
responsable de su lugar en la historia.

Esta conciencia no es un lujo teórico; es una exigencia que se impone
cuando la música se vive como un acto real en el tiempo y en el
espacio: cuando se toca en territorios donde el Estado no hace
presencia efectiva, cuando se comparte con músicos que aprendieron a
escuchar antes que a leer partituras. Esta dimensión política se
vuelve explícita en esas regiones donde la ausencia del Estado es
llenada por otros poderes; allí la música puede ser tanto un puente
como una amenaza, y el músico, sin buscarlo, se convierte en un
mediador entre mundos que no se hablan.

Pero mi posición no se agota en lo político. A mis 53 años, me
considero un estudiante permanente. Mi formación incluye también las
ciencias de la computación, y soy un entusiasta convencido del
software libre. No del "código abierto" —término que diluye el
compromiso ético—, sino del software libre tal como lo defiende
Richard Stallman: una cuestión de libertad, no de precio. Aprecio la
filosofía de Unix, pero mi lealtad está con GNU y el kernel Linux.
Mis herramientas cotidianas son Emacs/Lisp, la shell Bash, R para el
análisis estadístico, y SuperCollider con sus ramificaciones
(SuperDirt, FoxDot) para la síntesis y la composición algorítmica. El
kernel que utilizo está compilado y optimizado para baja latencia de
audio: una decisión técnica que es también una decisión filosófica,
porque afirma que la materia sonora merece un tratamiento riguroso
también en el dominio digital.

Creo en una educación basada en el \textbf{trivium} (gramática, lógica,
retórica) y el \textbf{quadrivium} (aritmética, geometría, música,
astronomía): las siete artes liberales que estructuraron el
pensamiento medieval y que, en plena era del silicio, siguen siendo
un modelo de formación integral. Para mí es natural hacer puentes
entre la música y el cálculo lambda, entre la organología y la
arquitectura de un sistema operativo, entre la notación musical y un
lenguaje de programación. Estoy acostumbrado a los cambios de
paradigma (παράδειγμα, \textbf{parádeigma}: modelo, patrón), como los
describe Kuhn en su epistemología.

Soy materialista, en el sentido preciso del materialismo filosófico
de Gustavo Bueno y su desarrollo musical por Vicente Chuliá. No soy
un repetidor de este sistema; observo que aporta herramientas
poderosas para pensar la música y la academia, y las utilizo
críticamente. Mi posición se sostiene sobre los tres pilares del
pluralismo buenista: la \textbf{symploké} (συμπλοκή) —el entrelazamiento de
las cosas, donde ni todo está conectado con todo ni nada está
conectado con nada—, la discontinuidad de la materia —que impide
reducirla a una unidad abstracta— y el no-monismo —que rechaza tanto
el monismo idealista como el monismo materialista grosero—. Estos
tres pilares no son independientes: la \textbf{symploké} es la forma misma
del entrelazamiento entre materialidades discontinuas; la
discontinuidad es la condición ontológica que hace imposible el
monismo; y el no-monismo es la consecuencia de reconocer que la
materia se da en géneros irreductibles. Mi intención en la academia
no es obtener un título; es la investigación.

Como decía Gabriel García Márquez, hay dos caminos para el escritor:
bajar el nivel de su escritura hasta el analfabetismo funcional, o
escribir según su propio nivel, elevando la exigencia cultural y
esperando que algún día el nivel del pueblo pueda llegar a esa altura.
No es clasismo: es apostar por la inteligencia del lector. Esa es la
apuesta de este ensayo.

Mi trayectoria es, en ese sentido, un caso particular de una pregunta
universal: ¿cómo se construye conocimiento musical cuando el camino
académico se interrumpe, no por fracaso, sino por la fuerza de la
propia experiencia?
\section{Trayectoria Formativa: Dialéctica entre Academia y Experiencia}
\label{sec:orgf14250b}
\subsection{Formación académica formal}
\label{sec:org67a11b4}
\subsubsection{Orquesta Sinfónica Juvenil de Bogotá}
\label{sec:orge519e66}
A los quince años inicié mis estudios básicos en la Orquesta
Sinfónica Juvenil de Bogotá, donde adquirí las bases de la
interpretación musical en formato orquestal. Mi instrumento principal
fue la guitarra clásica, y tuve como maestro a Andrés Sánchez,
reconocido contrabajista de la Orquesta Sinfónica de Colombia,
arreglista y guitarrista. Con él tomé mis primeros rudimentos de
solfeo y senté las bases técnicas que me acompañarían durante toda mi
trayectoria.
\subsubsection{Servicio militar — Banda de Guerra}
\label{sec:org3fa6879}
Durante el servicio militar obligatorio, me desempeñé como cornetista
y trompetista en la Banda de Guerra. Allí aprendí a ejecutar el "toque
de silencio": una obra musical simple, de pocas notas, que se toca con
corneta o trompeta para rendir homenaje a los soldados caídos en
combate. En total fueron doce toques de silencio. Pero esa es solo una
cara de la moneda: nunca me preguntaron si debía tocarse también
cuando muere un guerrillero o un paramilitar. Esa pregunta me acompañó
durante años, y aún no tiene respuesta.
\subsubsection{Universidad INCA de Colombia}
\label{sec:org79af5c7}
Estudié tres años en la Universidad INCA de Colombia: dos años como
guitarrista clásico y un tercer año en el programa de jazz, donde
aprendí de los ensambles que dirigía el músico Edy Martínez.
\subsubsection{Universidad de Cundinamarca — Facultad de Música}
\label{sec:orga670940}
Tuve la oportunidad de estudiar en la Universidad de Cundinamarca,
Facultad de Música, con una beca obtenida por haber sido el mejor
estudiante en el primer semestre. Allí conocí a Carlos Julio Ramírez,
Premio Nacional de Composición en Colombia. Él me alentó a estudiar
composición y me instó a profundizar en matemáticas y física para
entender las relaciones entre el número aritmético y el número
geométrico. De allí la música como número en el tiempo y en el espacio.

Fue él quien me puso en el centro del debate entre Pitágoras y
Filolao, y me ayudó a tomar posición: soy no monista. La materia es
discontinua; no se reduce a una unidad abstracta.

En la Facultad de Música entablé también una relación fundamental con
el maestro Jaime Cardona, decano de la facultad. Fue él quien me
relacionó con músicos de la talla del maestro Mario Posada, primer
violinista de la Orquesta Sinfónica de Colombia. El maestro Cardona
era un puente viviente hacia la gran tradición sinfónica del país.

El maestro Mario Posada me dio los rudimentos del violín. Con él
aprendí no solo técnica, sino una forma de escuchar: la precisión del
arco, la colocación del vibrato, la disciplina del ensayo.

No toco varios instrumentos por un deseo profundo: meterme en los
zapatos del músico que los toca, comprender sus capacidades y
límites, y así poder componer mejor. La guitarra es mi instrumento
principal; el piano es mi herramienta pedagógica y el ordenador mi
DAW.
\subsection{Formación autodidacta en Ciencias de la Computación y Software Libre}
\label{sec:org78bd9b6}
De manera paralela a mi trayectoria musical, y por iniciativa propia,
he cultivado una formación rigurosa en ciencias de la computación. No
fue una imposición curricular: fue un acto de libertad. Estudié los
materiales de \textbf{Open Source University} —repositorio en GitHub que
compila currículos completos de ciencias de la computación usando
exclusivamente recursos libres—, y mantengo una práctica constante de
programación en entornos libres.

Esta formación no es un añadido extracurricular: es parte integral de
mi pensamiento musical y filosófico. La relación entre música y
computación no es metafórica ni es una identidad; es un caso de
\textbf{symploké}: hay entrelazamiento sin reducción. El cálculo lambda
(\(\lambda\)) de Alonzo Church, base teórica de los lenguajes de
programación funcional, comparte con la música una estructura de
aplicación, abstracción y recursividad.
\subsubsection{Emacs/Lisp como entorno filosófico}
\label{sec:org97f62a0}
Utilizo Emacs como entorno integrado: allí escribo partituras en
LilyPond, documento en org-mode, programo en Elisp, analizo datos en
R, controlo SuperCollider para síntesis de audio en tiempo real, y
ejecuto scripts en Bash. Tengo funciones Elisp para casi todo. Esta
elección no es técnica: es filosófica. Emacs no es un editor de texto;
es un sistema operativo construido sobre Lisp. Cada función que escribo
en Elisp es un experimento ontológico: el ciclo READ-EVAL-PRINT de
Lisp es el mismo ciclo de la improvisación musical: propuesta,
ejecución, escucha, corrección.
\subsubsection{Kernel real-time y JackControl: la materia del tiempo}
\label{sec:orge35b871}
El kernel de mi sistema está compilado con opciones de \textbf{preemption} y
baja latencia (\textbf{real-time}), porque sé —materialistamente— que el
sonido es materia en el tiempo. Un \textbf{buffer underrun} no es un error
técnico; es una interrupción de la continuidad material de la obra.
JackControl es mi herramienta para hackear la tarjeta de sonido:
enrutar audio entre aplicaciones, conectar SuperCollider con Hydrogen,
enviar la guitarra eléctrica a través de efectos mientras el código
sigue sonando. A veces el portátil colapsa —un xrun, un buffer
underrun—, pero el músico sigue adelante. No hay errores en la música,
decía Miles Davis: se generan oportunidades.
\subsubsection{El método: cambiar pantallas, heredar del jazz}
\label{sec:org3f56a86}
Mi metodología de trabajo es una herencia del jazz: cambio de
pantallas para ir viendo el código preestablecido, lo modifico en la
marcha, lo llevo al extremo. No parto de la nada —*ex nihilo*—:
modifico, forkéo, reapropio. La mayoría de mis códigos son
modificaciones de códigos libres. Como dicta la ética hacker: la
información debe ser libre, el conocimiento no tiene dueño, y el
mérito se juzga por la calidad del trabajo, no por los títulos.
\subsection{Formación empírica y experiencia de campo}
\label{sec:org839fc6f}
\subsubsection{Territorios sin presencia efectiva del Estado: el laboratorio ontológico}
\label{sec:org45d20f9}
He tocado en territorios del Meta, en la región del Ariari, donde la
presencia efectiva del Estado colombiano ha sido históricamente
intermitente o nula debido al conflicto armado interno. No los llamo
"zonas rojas": ese es un eufemismo que despolitiza una realidad
geopolítica compleja. Para ingresar a estos territorios debes ir como
misionero claretiano; pasas puestos de control de la policía y el
ejército; luego encuentras retenes del paramilitarismo; y finalmente
un último control: el guerrillero. Allí exigen una "audición" para
decidir si lo que se va a hacer (música) es políticamente aceptable.

Esto no es anecdótico; es real. Y en esos momentos uno se hace
preguntas que la academia no suele formular: ¿qué es la música para
la derecha? Fácil: un himno nacional bien tocado, un tiple bien
peinado, una ranchera que hable de traición pero no de sindicatos.
¿Qué es la música para la izquierda? Un himno también, pero con el
puño en alto. ¿Y para los de arriba? Un concierto de gala, un Mozart
en el Teatro Colón, un patrocinio corporativo. ¿Y para los de abajo?
Un vallenato en la tienda, un rap que denuncie los "falsos positivos".
Y luego está uno, en medio de todo eso, con una guitarra al hombro y
un "toque de silencio" que no distingue entre muertos.
\subsubsection{Trabajo con músicos naturales}
\label{sec:orgc1bfeef}
He trabajado con músicos naturales: aquellos que aprendieron a tocar
de oído, sin partituras. Allí aprendí que el conocimiento musical no
se impone, sino que se entrelaza (\textbf{symploké}). Mi tarea no era
enseñarles, sino encontrar estrategias para que ellos tuvieran
herramientas sin perder su esencia.
\subsubsection{Distrito de Aguablanca, Cali — Producción musical en la periferia}
\label{sec:org7688642}
He estado en guetos como el Distrito de Aguablanca en Cali. Allí las
personas con poco acceso a la educación quieren expresarse como
pueden, y encuentran en la música de las periferias —como el hip hop
crítico de Asilo 38 y el álbum *Anarcolombia*— un grito contra una
sociedad violenta. En ese álbum participé como arreglista y productor.
He visto morir a músicos con talento inmenso; compañeros que cayeron
en el camino. He entendido que la música en Colombia no es un lujo:
es un acto de supervivencia, una forma de resistencia y de memoria.
\subsubsection{La Tropa Malagamba: una biografía de \textbf{symploké}}
\label{sec:org2dc8399}
A principios de los años dos mil, en una montaña de La Calera, conocí
a Francisco Magyarof —Panchis—, guitarrista gitano formado en la
tradición de Django Reinhardt. Desde el primer acorde, hubo
\textbf{symploké}: dos tradiciones discontinuas —el jazz gitano y la
formación clásica— que se entrelazaban sin fundirse.

Panchis tenía el flamenco interiorizado, pero lo interpretaba de una
manera que los académicos llamarían "sucia". No era un defecto: era
un estilo. Con él aprendí los palos flamencos. Todo era oral:
imitación, admiración mutua. Yo escribía en notación musical lo que
tocábamos —\(M_1\) sobre el papel, \(M_3\) en gestación—. Fue entonces
cuando descubrí la progresión napolitana y entendí, por fin, qué era
el fraseo.

Con él fundé la Tropa Malagamba, un grupo callejero que incorporaba
bandola llanera, tiple, viola, banjo y guitarra. Luego llegó César
Reyes, sobrino del maestro Ernesto Díaz —fundador de la Orquesta
Sinfónica de Colombia—, tocando la viola. El grupo creció hasta
convertirse en Ganjariquies: una banda eléctrica de quince
integrantes, parte del movimiento de Resistencia Cultural que, con
hip hop, ska y reggae, denunciaba los crímenes de Estado en pueblos
donde la presencia del gobierno era un rumor.

Para entrar a esos territorios había que ir como misionero
claretiano. Los Misioneros Claretianos fueron nuestro salvoconducto.

César Reyes —el violista de la Tropa— era sobrino del maestro Ernesto
Díaz. Yo estudié en esa orquesta cuando era adolescente. César me
presentó a su madre, Teresa, sobrina del maestro Ernesto. Y a su
padrastro: Andrés Sánchez. Mi primer maestro de guitarra. El hombre
que me enseñó los rudimentos del solfeo cuando yo tenía quince años.
La \textbf{symploké} no es un concepto. Es una biografía.

Las categorías del materialismo filosófico no eran teoría: eran
práctica. La materia sonora (\(M_1\)) vibraba entre la guitarra gitana
y la barroca; la percepción (\(M_2\)) de un público que bailaba sin
saber qué era una progresión napolitana; y el concepto (\(M_3\)) se
materializaba en letras que hoy son documentos históricos.

Hace unos años perdí el contacto con mi grupo. La música en Colombia
no es entretenimiento: es supervivencia. Y a veces, ni siquiera eso
alcanza.
\section{Marco Teórico: Filosofía Materialista de la Música}
\label{sec:org5638919}
\subsection{Los Tres Pilares del Pluralismo de Gustavo Bueno}
\label{sec:orgd44dce5}
Antes de desarrollar los géneros de materialidad, es necesario
explicitar los tres pilares sobre los que se sostiene el pluralismo
materialista de Gustavo Bueno, pues son el fundamento de todo el
edificio teórico que aquí se despliega. No se trata de tres ideas
independientes, sino de tres aspectos de una misma ontología
(ὀντολογία, \textbf{ontología}: discurso sobre el ente) que se implican
mutuamente (??, ????) (??, a).
\subsubsection{Symploké (συμπλοκή): el entrelazamiento de las cosas}
\label{sec:org533e19c}
El término \textbf{symploké} proviene de Platón y significa
"entrelazamiento". Bueno lo recupera para designar la posición
ontológica que se sitúa entre dos extremos igualmente erróneos: el
monismo (todo está conectado con todo en una unidad absoluta) y el
atomismo radical (nada está conectado con nada, solo hay partículas
aisladas). La \textbf{symploké} afirma que unas cosas están conectadas con
otras, pero no todas con todas. Hay conexiones reales, pero también
hay desconexiones reales.

\begin{verbatim}
library(ggplot2)

crear_circulo <- function(x0, y0, r, n = 200) {
  theta <- seq(0, 2*pi, length.out = n)
  data.frame(x = x0 + r * cos(theta), y = y0 + r * sin(theta))
}

c1 <- crear_circulo(0, 0.5, 1.3)
c1$pilar <- "Symploké"
c2 <- crear_circulo(-0.9, -0.5, 1.3)
c2$pilar <- "Discontinuidad\nde la materia"
c3 <- crear_circulo(0.9, -0.5, 1.3)
c3$pilar <- "No-monismo"

circulos <- rbind(c1, c2, c3)

etiquetas <- data.frame(
  x = c(0, -1.5, 1.5),
  y = c(1.2, -1.5, -1.5),
  label = c("Symploké\n(συμπλοκή)", "Discontinuidad\nde la materia", "No-monismo\n(μόνος)")
)

ggplot() +
  geom_polygon(data = circulos, aes(x = x, y = y, fill = pilar, group = pilar), 
               alpha = 0.35) +
  geom_text(data = etiquetas, aes(x = x, y = y, label = label), 
            size = 5, fontface = "bold", color = "gray20") +
  annotate("text", x = 0, y = -2.2, 
           label = "Materialismo Filosófico — Escuela de Toledo", 
           size = 5.5, fontface = "italic", color = "gray10") +
  annotate("text", x = 0, y = 2.5, 
           label = "Los Tres Pilares del Pluralismo de Gustavo Bueno", 
           size = 6, fontface = "bold") +
  scale_fill_manual(values = c("#2B579A", "#B7472A", "#D4A017"), guide = "none") +
  coord_fixed() +
  theme_void()
\end{verbatim}

\begin{figure}[htbp]
\centering
\includegraphics[width=.9\linewidth]{/home/spock/Documents/symploke-sonora/img/tres_pilares.png}
\caption{Los Tres Pilares del Pluralismo de Gustavo Bueno}
\end{figure}
\subsubsection{Discontinuidad de la materia}
\label{sec:org57beefd}
La materia no es una sustancia (\textbf{substantia}, traducción latina del
griego ὑπόστασις, \textbf{hypóstasis}: lo que subyace) unitaria que se
despliega en diversos aspectos, sino que se da en géneros de
materialidad irreductibles entre sí. \(M_1\) (materia física), \(M_2\)
(materia perceptiva) y \(M_3\) (materia conceptual) no son grados de una
misma realidad, sino discontinuidades ontológicas.
\subsubsection{No-monismo}
\label{sec:org8b1a213}
Si la materia es discontinua, entonces no puede ser reducida a un
principio único (μόνος, \textbf{mónos}: único, solo). El no-monismo rechaza
tanto el monismo idealista como el monismo materialista grosero.
Frente a ambos, afirma la pluralidad irreductible de los géneros de
materialidad. La música no es solo vibración (\(M_1\)), ni solo
percepción (\(M_2\)), ni solo concepto (\(M_3\)): es las tres, pero
precisamente porque ninguna de ellas puede absorber a las otras dos.
\subsection{Antecedentes Filosóficos: De los presocráticos a Aristóteles}
\label{sec:org68b031c}
El materialismo filosófico no surge de la nada. Sus tesis —la
discontinuidad de la materia, el entrelazamiento de géneros, la música
como ontología— fueron anticipadas, negadas o desviadas por la
tradición filosófica anterior. Recuperar esa tradición no es un
ejercicio de arqueología: es un acto de justicia epistémica.

Pero antes de la tradición escrita, hubo música. Los hallazgos
arqueológicos lo confirman: en la cueva de Hohle Fels, en el suroeste
de Alemania, se encontraron flautas talladas en hueso de buitre y
marfil de mamut datadas entre 40,000 y 43,000 años de antigüedad
(??, a). Son los instrumentos musicales más antiguos
conocidos fabricados por \textbf{Homo sapiens}. En Divje Babe, Eslovenia, un
fémur de oso de las cavernas con perforaciones —datado entre 43,000 y
60,000 años— podría ser una flauta neandertal, aunque su origen sigue
siendo debatido (??, a).

Estos hallazgos revelan algo fundamental: la música no es un invento
de la civilización. Es una práctica material que antecede a la
escritura, a la agricultura, al Estado. Antes de que existiera la
palabra «filosofía», ya existía la vibración organizada. El sonido no
necesita concepto (\(M_3\)) para existir: necesita materia que vibre
(\(M_1\)) y un oído que la perciba (\(M_2\)).

Los filósofos griegos no inventaron la música. Heredaron una práctica
milenaria y la sometieron a reflexión sistemática. Lo que sigue es el
rastreo de esa reflexión, desde los presocráticos hasta Aristóteles,
atendiendo no solo a lo que pensaron, sino a las tecnologías concretas
que tenían a su disposición para pensarla.
\subsubsection{Los presocráticos: el número, el devenir y el ser}
\label{sec:org141579d}
La filosofía nace con una pregunta que sigue siendo la nuestra: ¿de
qué está hecho el mundo? Los presocráticos respondieron de maneras que
aún resuenan en la ontología musical.

Tales de Mileto (c. 624-546 a.C.) propuso que el principio de todas
las cosas es el agua. No es una afirmación ingenua: es la primera tesis
materialista de la historia. La materia no es creada por los dioses;
es eterna y se transforma. Anaximandro (c. 610-546 a.C.) fue más
preciso: el principio no puede ser un elemento concreto, sino algo
indeterminado, el ἄπειρον (\textbf{ápeiron}: lo ilimitado). Esta intuición
—que el origen de la materia no puede ser una sustancia particular—
reaparecerá en Filolao y, dos milenios después, en el materialismo
filosófico: la materia no se reduce a un principio único.

Anaxímenes (c. 585-524 a.C.) identificó el aire como principio y
explicó la transformación por condensación y rarefacción. La
tecnología de su época era la observación directa: el aire se
condensa en agua, el agua en tierra. No había instrumentos de
medición, pero sí una intuición materialista: la materia se transforma
sin necesidad de dioses.

Heráclito de Éfeso (c. 535-475 a.C.) introdujo una idea que atraviesa
todo este ensayo: el devenir. «No nos bañamos dos veces en el mismo
río», escribió. La realidad no es estática; es proceso. Y el proceso
no es caos: está regulado por el λόγος (\textbf{lógos}), una proporción, una
medida. La música es la encarnación más perfecta de esa proporción:
sonidos que devienen en el tiempo según relaciones matemáticas.

Parménides de Elea (c. 515-450 a.C.) tomó el camino opuesto: el ser
es uno, inmóvil, eterno. El cambio es una ilusión. Esta tensión entre
Heráclito y Parménides —devenir contra permanencia, pluralidad contra
unidad— es la matriz de todas las discusiones ontológicas posteriores.
Y es también la tensión que recorre la historia de la música: ¿es la
música un flujo que nunca se repite, o es una estructura que subyace
bajo las apariencias sonoras?
\subsubsection{Pitágoras y el monocordio: la tecnología al servicio de la ontología}
\label{sec:orgca8ac2a}
Pitágoras de Samos (c. 570-495 a.C.) merece un lugar aparte. No solo
fue el primero en matematizar la música, sino que lo hizo con una
tecnología concreta: el monocordio (μονόχορδον, \textbf{monóchordon}: una
sola cuerda). Este instrumento, descrito por Boecio en \textbf{De
institutione musica} (Libro I, capítulos 10-11) y por Euclides en los
\textbf{Elementos de armonía}, consistía en una cuerda tensada sobre una caja
de resonancia con un puente móvil que permitía dividirla en
proporciones exactas. No era un instrumento musical en sentido
estricto: era un instrumento científico. Un laboratorio de física
acústica portátil.

Al dividir la cuerda en mitades (2:1), Pitágoras obtenía la octava; en
tercios (3:2), la quinta; en cuartos (4:3), la cuarta. El monocordio
no era un instrumento para hacer música: era un instrumento para hacer
ontología. La tecnología no era un medio para producir sonido, sino
para producir conocimiento.

\begin{verbatim}
library(ggplot2)

monocordio <- data.frame(
  proporcion = factor(c("1:1", "2:1", "3:2", "4:3", "5:4", "6:5"),
                      levels = c("1:1", "2:1", "3:2", "4:3", "5:4", "6:5")),
  ratio = c(1, 2, 1.5, 1.333, 1.25, 1.2),
  intervalo = c("Unísono", "Octava", "Quinta", "Cuarta", "Tercera Mayor", "Tercera menor"),
  descubridor = c("Unísono", "Pitágoras", "Pitágoras", "Pitágoras", "Posterior", "Posterior")
)

ggplot(monocordio, aes(x = proporcion, y = ratio, fill = descubridor)) +
  geom_col(width = 0.6, color = "white", linewidth = 0.5) +
  geom_text(aes(label = intervalo), vjust = -0.5, size = 4.5, fontface = "italic") +
  scale_fill_manual(values = c("Pitágoras" = "#D4AC0D", 
                                "Posterior" = "#922B21",
                                "Unísono" = "gray70"),
                    guide = "none") +
  geom_hline(yintercept = 1, linetype = "dashed", color = "gray50") +
  labs(title = "El Monocordio de Pitágoras",
       subtitle = "Proporciones matemáticas de las consonancias — s. VI a.C.",
       x = "Proporción de la cuerda",
       y = "Ratio de frecuencia") +
  theme_minimal(base_size = 14) +
  theme(plot.title = element_text(face = "bold", size = 18),
        plot.subtitle = element_text(color = "gray40"))
\end{verbatim}

\begin{figure}[htbp]
\centering
\includegraphics[width=.9\linewidth]{/home/spock/Documents/symploke-sonora/img/monocordio_pitagoras.png}
\caption{Proporciones del monocordio pitagórico}
\end{figure}

Veintiséis siglos después, el mismo experimento puede replicarse con
SuperCollider:

\begin{verbatim}
// Monocordio pitagórico en SuperCollider
(
SynthDef(\monocordio, { |out=0, freq=440, amp=0.3, pan=0|
    var sig = SinOsc.ar(freq, 0, amp) * EnvGen.kr(
        Env.perc(0.01, 1.5, 1, -4),
        doneAction: 2
    );
    Out.ar(out, Pan2.ar(sig, pan));
}).add;

Routine({
    [1, 2, 1.5, 4/3, 1.25, 1.2, 1].do { |ratio, i|
        Synth(\monocordio, [\freq, 440 * ratio]);
        (1.5).wait;
    };
}).play;
)
\end{verbatim}

Pero esta matematización tuvo un precio. Pitágoras convirtió el número
en el principio de todas las cosas. La música no era sonido: era
proporción. El cosmos entero, según la escuela pitagórica, era una
armonía que el oído humano no podía escuchar pero que el intelecto
podía comprender. Las esferas celestes, al moverse, producían una
música que solo los sabios podían percibir.

Aquí anida el idealismo pitagórico. Porque las ondas sinusoidales,
hoy lo sabemos, no se propagan en el vacío del espacio. No hay sonido
en el cosmos: hay vibraciones electromagnéticas, campos, radiación de
fondo. Pero no hay música de las esferas. Lo que Pitágoras tomó por
armonía celeste era, en realidad, una proyección de su propia
tecnología: el monocordio no medía el universo; medía una cuerda
tensada en una caja de madera. Confundir el instrumento con el cosmos
es el pecado original del idealismo musical.
\subsubsection{Arquitas de Tarento y la acústica materialista}
\label{sec:orgd628716}
Arquitas de Tarento (c. 428-347 a.C.), contemporáneo de Platón, fue
el primero en formular una teoría física del sonido que no recurría a
la armonía celeste. El sonido, según Arquitas, es vibración que se
transmite por el aire. No es magia, no es número puro: es mecánica.
Construyó autómatas —incluyendo una paloma de madera que volaba con
vapor— y se le atribuye la invención de la polea y el tornillo. Su
enfoque tecnológico influyó en la acústica griega y anticipó, veinte
siglos antes, la física del sonido como fenómeno material.
\subsubsection{Platón: la música como educación del alma}
\label{sec:org065cbe9}
Platón (c. 427-347 a.C.) hereda el pitagorismo pero lo lleva al
terreno político y pedagógico. En \textbf{La República}, la música no es un
entretenimiento: es un instrumento de formación del ciudadano. Los
modos musicales no son neutros: cada uno produce un efecto moral
distinto. El modo dorio infunde valor; el modo lidio, molicie. Por eso
la música debe ser regulada por el Estado.

Desde el materialismo filosófico, esta concepción es un avance y un
retroceso. Avance: Platón reconoce que la música tiene efectos reales
sobre la materia perceptiva (\(M_2\)). Retroceso: subsume la música a la
política y la moral, negándole autonomía ontológica. La música no es
valiosa por sí misma, sino por sus efectos. Es un medio, no un fin.
\subsubsection{Aristóteles y las Categorías (Κατηγορίαι, \textbf{Kategoríai})}
\label{sec:org16587f1}
Aristóteles (384-322 a.C.) es el primer filósofo que ofrece un sistema
de categorías para describir la realidad sin reducirla a un principio
único. Frente al monismo pitagórico y al idealismo platónico,
Aristóteles propone una ontología pluralista: el ser se dice de muchas
maneras.

En las \textbf{Categorías}, Aristóteles distingue diez modos de predicar el
ser: sustancia (οὐσία, \textbf{ousía}), cantidad (ποσόν, \textbf{posón}), cualidad
(ποιόν, \textbf{poión}), relación (πρός τι, \textbf{prós ti}), lugar (ποῦ, \textbf{poû}),
tiempo (ποτέ, \textbf{poté}), posición (κεῖσθαι, \textbf{keîsthai}), posesión
(ἔχειν, \textbf{ékhein}), acción (ποιεῖν, \textbf{poieîn}) y pasión (πάσχειν,
\textbf{páskhein}).

Aplicadas a la música, estas categorías permiten afirmar que la música
no es una sustancia abstracta, sino un ente que se da en el tiempo y
en el espacio; que tiene cantidad (duración, altura), cualidad (timbre,
intensidad), relación (intervalos, armonía), lugar (la sala, el
instrumento), tiempo (el tempo, la agógica) y posición (el gesto del
director sobre el espacio eufónico). La música no está fuera de las
categorías; está en ellas.

Pero las categorías mismas son discontinuas entre sí: la cantidad no se
reduce a la cualidad, ni la relación al lugar. La música no es una
unidad; es un entrelazamiento (\textbf{symploké}) de determinaciones
categoriales irreductibles unas a otras (??, a).

La tecnología de Aristóteles no era el monocordio, sino el teatro
griego. En \textbf{Problemas} (Libro XI), analiza por qué el sonido se
propaga mejor en ciertas condiciones arquitectónicas. Las máscaras de
los actores tenían una función acústica, no solo dramática: amplificaban
la voz. La música no era solo teoría: era arquitectura. Esta atención
a las condiciones materiales del sonido hace de Aristóteles un
precursor del materialismo musical, aunque su teoría de la música como
catarsis lo mantuviera aún dentro del paradigma idealista.
\subsubsection{Del Liceo a Roma: los estoicos y la música como orden cósmico}
\label{sec:orgeb98147}
Los estoicos —Zenón de Citio (c. 334-262 a.C.), Séneca (4 a.C.-65
d.C.), Epicteto (55-135 d.C.)— heredaron la idea platónica de que la
música refleja el orden del cosmos, pero la desvincularon de la
política. Para los estoicos, la música es una expresión del λόγος
(\textbf{lógos}) universal, la razón que gobierna el universo. Escuchar
música es sintonizar con el orden cósmico.

Séneca, en sus \textbf{Cartas a Lucilio}, escribió que la música es una de
las artes liberales que preparan el alma para la virtud. Pero advirtió
contra el exceso: la música no debe convertirse en un fin en sí misma,
sino en un medio para la sabiduría. Esta tensión entre la música como
arte liberal y la música como mero entretenimiento recorrerá toda la
Edad Media hasta llegar a Boecio y San Isidoro.
\subsubsection{De Roma a la Ilustración: un puente necesario}
\label{sec:orgef83d9a}
Entre Séneca y Wolff transcurren doce siglos. No es posible —ni
necesario— recorrerlos todos. Pero sí es necesario trazar un puente
que conecte la música como orden cósmico (los estoicos) con la música
como objeto de la estética (Baumgarten). Ese puente tiene hitos:
filósofos, matemáticos, músicos y científicos que, con las tecnologías
de su tiempo, prepararon el terreno para que la música se convirtiera
en un problema filosófico autónomo.
\begin{enumerate}
\item Boecio (c. 480-524 d.C.): las tres músicas
\label{sec:org35f6a55}
Boecio, en \textbf{De institutione musica}, recupera el monocordio pitagórico
y establece una clasificación que dominará todo el pensamiento
medieval: la música \textbf{mundana} (la armonía del cosmos), la música
\textbf{humana} (la armonía del cuerpo y el alma) y la música \textbf{instrumentalis}
(la que producen los instrumentos). La tecnología del monocordio
—una cuerda, un puente móvil, una caja de resonancia— permitía
demostrar empíricamente las proporciones que los pitagóricos habían
intuido. Boecio no era un músico: era un matemático que usaba el
sonido para enseñar aritmética.
\item San Isidoro de Sevilla (c. 560-636): el \textbf{quadrivium}
\label{sec:org29741da}
San Isidoro, en las \textbf{Etimologías}, consagra la música como una de las
cuatro artes matemáticas del \textbf{quadrivium}: aritmética, geometría,
música y astronomía. Su frase —«sin la música, ninguna disciplina
puede ser perfecta»— no es un elogio: es una definición. La música no
es adorno: es fundamento.
\item Al-Farabi (872-950) y Avicena (980-1037): el puente islámico
\label{sec:org66b97d1}
Mientras Europa entraba en la Edad Media, el mundo islámico preservaba
y expandía la tradición griega. Al-Farabi, en el \textbf{Gran libro de la
música}, sistematizó la teoría musical árabe y la conectó con la
física aristotélica. Avicena, en el \textbf{Canon de la medicina}, estudió
los efectos terapéuticos del sonido sobre el cuerpo. La música no era
solo matemática: era medicina.
\item Guido de Arezzo (c. 991-1033): la notación como tecnología
\label{sec:orgc1d27ec}
Guido de Arezzo no fue un filósofo, pero su invención —la notación
musical en tetragrama— transformó la música en un objeto de estudio.
Antes de Guido, la música se transmitía oralmente. Después de Guido,
se podía escribir, leer, comparar, analizar. La notación es una
tecnología de fijación que permite la acumulación de conocimiento.
Sin Guido, no hay tratado. Sin tratado, no hay ciencia musical.
\item Gioseffo Zarlino (1517-1590): la armonía como ciencia
\label{sec:orgc185570}
En el Renacimiento, Zarlino sistematizó el contrapunto en \textbf{Le
institutioni harmoniche} (1558). Su tecnología era la imprenta: por
primera vez, las partituras y los tratados se difundían masivamente.
Zarlino describió los acordes triádicos y sus inversiones, sentando
las bases de la armonía moderna. La música dejaba de ser un arte de
oficio para convertirse en una ciencia de proporciones.
\item Vincenzo Galilei (c. 1520-1591): la crítica al pitagorismo
\label{sec:orgf229449}
Vincenzo Galilei, padre de Galileo, fue el primero en criticar
empíricamente el pitagorismo musical. En su \textbf{Diálogo de la música
antigua y moderna} (1581), demostró que la relación entre la tensión
de una cuerda y su frecuencia no seguía una ley lineal, sino que
dependía de factores materiales: el grosor, el material, la densidad.
La música no es número puro (\(M_3\)): es materia que vibra (\(M_1\)).
Su tecnología era el laúd y la experimentación doméstica. Su
heredero sería Galileo, que aplicaría el mismo método al cosmos.
\item Johannes Kepler (1571-1630): el último pitagórico
\label{sec:orgdd5bd7c}
Kepler, en \textbf{Harmonices Mundi} (1619), intentó demostrar que las
órbitas planetarias correspondían a intervalos musicales. Su
tecnología era el telescopio: sin él, no habría podido observar las
elipses planetarias. Kepler fue el último gran intento de unificar
astronomía y música. Su fracaso —las órbitas no son círculos
perfectos, sino elipses— es una lección ontológica: el cosmos no es
una sinfonía. Pero su intuición —que las mismas leyes matemáticas
gobiernan las órbitas y las cuerdas— era correcta.
\item René Descartes (1596-1650): la música como mecánica
\label{sec:org5dffe50}
Descartes, en su \textbf{Compendium Musicae} (1618), trató la música como un
estímulo mecánico de las pasiones. El sonido golpea el oído, el oído
transmite la vibración al cerebro, el cerebro produce una emoción.
Todo es mecánica. La música no es un lenguaje del alma: es una cadena
causal de eventos físicos. Esta reducción mecanicista es,
paradójicamente, un avance materialista: la música no necesita
espíritu para funcionar.
\item Isaac Newton (1643-1727): el cálculo como lenguaje
\label{sec:orgcff7f3d}
Newton no escribió sobre música. Pero su cálculo infinitesimal
—desarrollado en paralelo con Leibniz— proporcionó la herramienta
matemática que permitiría analizar el sonido como un fenómeno
dinámico. La vibración de una cuerda, la propagación de una onda, la
descomposición de un timbre en armónicos: todo eso requiere cálculo.
Sin Newton, no hay Fourier. Sin Fourier, no hay análisis espectral.
Sin análisis espectral, no hay síntesis digital. La tecnología de
Newton —el cálculo— es, tres siglos después, la tecnología de
SuperCollider.
\item Gottfried Wilhelm Leibniz (1646-1716): la aritmética inconsciente
\label{sec:orged19003}
Leibniz, rival de Newton en la invención del cálculo, sí escribió
sobre música. Su definición —«la música es un ejercicio oculto de
aritmética del alma que no sabe que está contando»— es el puente
perfecto hacia Wolff y Baumgarten. La música es inconsciente, pero
es racional. El alma cuenta sin saberlo. Esta idea —que la percepción
musical es un proceso cognitivo, no solo sensorial— será el
fundamento de la estética como disciplina filosófica.
\item Christian Wolff (1679-1754): la sistematización
\label{sec:org47019ba}
Wolff, discípulo de Leibniz, fue el gran sistematizador de la
filosofía alemana anterior a Kant. Dividió el conocimiento en ciencias
racionales (lógica, metafísica) y ciencias empíricas (física,
psicología). La música quedó en una zona ambigua: ¿es racional o
empírica? ¿Es matemática o sensación? Wolff no resolvió la cuestión,
pero la formuló con una claridad que hizo posible la estética de
Baumgarten. Su tecnología era la imprenta y el aula universitaria: la
filosofía se enseñaba, se discutía, se sistematizaba.
\end{enumerate}
\subsubsection{Christian Wolff, Alexander Baumgarten y el nacimiento de la Estética}
\label{sec:org3508871}
El puente que conduce de Leibniz a Kant pasa por dos figuras
fundamentales: Christian Wolff (1679-1754) y Alexander Gottlieb
Baumgarten (1714-1762). Wolff sistematizó la filosofía alemana;
Baumgarten fundó la estética como disciplina autónoma. Ambos lo
hicieron con la tecnología de su tiempo: la imprenta, el aula
universitaria, el tratado sistemático.
\begin{enumerate}
\item Wolff y la división del conocimiento
\label{sec:org6b45188}
Wolff, discípulo de Leibniz, fue el gran sistematizador de la
filosofía alemana anterior a Kant. Dividió el conocimiento en ciencias
racionales (lógica, metafísica) y ciencias empíricas (física,
psicología). La música quedó en una zona ambigua: ¿es racional o
empírica? ¿Es matemática o sensación? Wolff no resolvió la cuestión,
pero la formuló con una claridad que hizo posible la estética.

Hay que entender el contexto teológico de esta división. Wolff era
protestante, formado en la escolástica luterana. El protestantismo
había despojado a los templos de imágenes y había reducido la música
litúrgica a lo esencial: el coral, la palabra cantada, el órgano como
acompañamiento. La belleza sensible —el oro, los frescos, la
polifonía compleja— era sospechosa de idolatría. La música no debía
ser un fin en sí misma: debía servir a la palabra divina.

El catolicismo, en cambio, había abrazado el Barroco como una
tecnología de persuasión sensorial. Las iglesias católicas del siglo
XVII eran máquinas de belleza: pinturas, esculturas, retablos dorados,
música polifónica que envolvía al fiel en una experiencia total. La
música no era un medio para la palabra: era un fin en sí misma. Bach,
protestante, compuso para la gloria de Dios. Vivaldi, católico,
compuso para la gloria del sonido.

Esta tensión entre catolicismo y protestantismo no es un detalle
teológico: es el marco material donde nace la estética. Porque la
estética de Baumgarten no es católica: es protestante. No busca la
belleza como esplendor sensorial, sino el conocimiento sensible como
forma inferior pero legítima de conocimiento.
\item Baumgarten y la invención de la Estética
\label{sec:orga3fcbc8}
Alexander Gottlieb Baumgarten acuñó el término "estética" (αἰσθητική,
\textbf{aisthetiké}: ciencia del conocimiento sensible) en su obra
\textbf{Aesthetica} (1750). Con este gesto, Baumgarten separó las bellas
artes de las artes liberales y las convirtió en un objeto de estudio
filosófico autónomo.

Antes de Baumgarten, las artes estaban dispersas. La música pertenecía
al \textbf{quadrivium} (junto con la aritmética, la geometría y la
astronomía). La pintura y la escultura eran oficios mecánicos. La
poesía era retórica. Baumgarten unificó estas disciplinas bajo un
nuevo concepto: las bellas artes, cuyo fin no es la utilidad ni la
verdad, sino la belleza.

Pero Baumgarten no definió la estética como una ciencia de la belleza,
sino como una ciencia del conocimiento sensible. La diferencia es
crucial. La belleza era un concepto platónico, objetivo, metafísico.
El conocimiento sensible era un proceso psicológico, subjetivo,
analizable. Baumgarten no preguntó "¿qué es bello?", sino "¿cómo
percibimos lo bello?". Es un giro copernicano: la estética no estudia
el objeto bello, sino el sujeto que percibe.

Desde el materialismo filosófico, este giro es ambiguo. Por un lado,
Baumgarten reconoce la legitimidad del conocimiento sensible (\(M_2\)),
que la tradición racionalista había despreciado. Por otro lado,
inaugura el subjetivismo estético: la belleza deja de ser una
propiedad de la materia (\(M_1\)) para convertirse en una experiencia
del sujeto (\(M_2\)). La estética idealista —que Kant llevará a su
culminación— comienza aquí.
\item La dialéctica entre estética y bellas artes
\label{sec:org1b87387}
La unificación de las bellas artes bajo el concepto de "estética" tuvo
una consecuencia paradójica: separó el arte de la ciencia. Antes de
Baumgarten, la música era una ciencia matemática (el \textbf{quadrivium}).
Después de Baumgarten, la música es un arte bello. La diferencia no es
terminológica: es ontológica.

Una ciencia se estudia con el intelecto. Un arte bello se experimenta
con el gusto. La ciencia busca la verdad. El arte busca la belleza. Al
separar la verdad de la belleza, Baumgarten consumó el divorcio entre
la música y las matemáticas que el \textbf{quadrivium} había mantenido unidas
durante un milenio.

Este divorcio es, en gran medida, el responsable de la situación que
este ensayo denuncia: la música se enseña en conservatorios, la
física en facultades de ciencias, y nadie le explica al músico que las
supercuerdas vibran ni al físico que la serie de armónicos es música.
\item Kant y el juicio estético
\label{sec:org1c41922}
Immanuel Kant (1724-1804) asumió el andamiaje de Wolff y Baumgarten y
lo transformó en la \textbf{Crítica del Juicio} (1790). Kant no era músico,
pero su teoría del juicio estético definiría la forma en que la
modernidad ha entendido la música durante más de dos siglos.

Para Kant, el juicio estético es desinteresado: no busca poseer el
objeto, ni utilizarlo, ni conocerlo. Solo busca el placer de la
contemplación. La música es un "juego de emociones" que afecta al
ánimo, pero que ofrece poca cultura intelectual en comparación con la
poesía. La música, según Kant, es más "goce auditivo" que "cultura",
y por eso la sitúa en un lugar bajo en su jerarquía de las artes.

Este veredicto ha pesado como una losa sobre la música occidental.
Kant consagra el idealismo estético que Baumgarten había inaugurado:
la música no es materia que vibra (\(M_1\)), ni percepción que organiza
(\(M_2\)), ni concepto que estructura (\(M_3\)). Es un "juego de
emociones". Un entretenimiento elevado, pero entretenimiento al fin.

Sin embargo, Kant también nos legó una herramienta que este ensayo
reivindica: el imperativo categórico. «Obra sólo según aquella máxima
por la cual puedas querer que al mismo tiempo se convierta en ley
universal». Traducido a los términos de este ensayo: si pido que mi
saber sea reconocido como válido, debo estar dispuesto a que
cualquier otro músico en mis mismas condiciones reciba el mismo
reconocimiento. No pido un privilegio: pido que se aplique una ley
que ya existe.

Kant, formado en la escolástica protestante, siguió un método que
debía mucho a Wolff: división sistemática, argumentación rigurosa,
desconfianza hacia la experiencia sensible como fuente de conocimiento.
Su tecnología era el libro impreso y la cátedra universitaria. Nunca
tocó un instrumento. Nunca asistió a un concierto del que se tenga
registro. Su conocimiento de la música era puramente conceptual
(\(M_3\)), sin anclaje en \(M_1\) ni en \(M_2\). Quizás por eso su juicio
sobre ella fue tan pobre.
\end{enumerate}
\subsubsection{Música y ciencia en el siglo XVII: dos tríos en paralelo}
\label{sec:orgf8a8e5f}
El siglo XVII fue un siglo de revoluciones. La ciencia abandonó a
Aristóteles; la música abandonó el Renacimiento. Dos tríos —uno de
científicos y filósofos, otro de músicos— transformaron la forma de
entender el mundo y el sonido. No se conocieron todos entre sí, pero
compartieron un mismo horizonte: la materia no es un cosmos cerrado,
sino un campo de fuerzas que puede ser medido, analizado y
transformado.
\begin{enumerate}
\item El trío científico: Galileo, Bacon, Descartes
\label{sec:org06fc1ea}
Galileo Galilei (1564-1642) apuntó el telescopio al cielo y demostró
que la materia celeste es igual a la terrestre. Pero Galileo también
era laudista: su padre, Vincenzo Galilei, fue el primer crítico
empírico del pitagorismo musical. La música de la familia Galilei no
era especulación: era experimentación.

Francis Bacon (1561-1626) no hizo música, pero su método inductivo
—observar, clasificar, experimentar— es el mismo que un músico aplica
cuando explora las posibilidades de un instrumento. El empirismo de
Bacon es, en el fondo, una ética de la atención: no imponer ideas a la
materia, sino dejar que la materia revele sus patrones.

René Descartes (1596-1650) escribió un tratado de música —el
\textbf{Compendium Musicae} (1618)— donde redujo la experiencia sonora a
mecánica corporal. El sonido golpea el oído, el oído transmite la
vibración al cerebro, el cerebro produce una emoción. Todo es cadena
causal. La música no es un misterio: es un mecanismo.
\item El trío musical: Monteverdi, Frescobaldi, Schütz
\label{sec:org28a06de}
Claudio Monteverdi (1567-1643) rompió con la polifonía renacentista y
fundó la \textbf{seconda pratica}: la música debe servir al texto, no al
revés. En su ópera \textbf{Orfeo} (1607), la voz humana es el centro y la
orquesta la envuelve. La tecnología de Monteverdi era la orquesta
barroca: violines, violas, bajos continuos, tiorbas. Un laboratorio de
timbres.

Girolamo Frescobaldi (1583-1643) fue el organista más influyente del
Barroco temprano. Sus \textbf{Fiori musicali} (1635) exploraban las
posibilidades del teclado como un científico explora un laboratorio:
cada registro, cada combinación de tubos, cada variación rítmica era
un experimento. Bach estudió sus obras un siglo después.

Heinrich Schütz (1585-1672) viajó a Venecia y absorbió el estilo de
Monteverdi. Lo llevó a Alemania, donde lo fusionó con la tradición
luterana del coral. Schütz es el eslabón que conecta el Renacimiento
italiano con el Barroco alemán. Sin Schütz, no hay Bach.
\item Paralelismos sin contacto directo
\label{sec:orgfca0ac7}
Monteverdi y Galileo compartieron ciudad (Venecia) y mecenas. Aunque
no hay registro de que se conocieran, el padre de Galileo —Vincenzo
Galilei— fue un teórico musical que desafió el pitagorismo con el
mismo espíritu con que Galileo desafió el aristotelismo. La música y
la ciencia compartían un mismo adversario: la autoridad incuestionada.

Frescobaldi y Bacon nunca se encontraron, pero su método era el mismo:
observación, variación, experimentación. El teclado de Frescobaldi era
su laboratorio.

Schütz y Descartes nunca se encontraron, pero compartieron una misma
convicción: la música es expresión de las pasiones. En Descartes, las
pasiones son mecánicas; en Schütz, son sonido. Pero ambos entendieron
que la música no es un ornamento: es una tecnología de transformación
afectiva.
\item La tecnología que los unió
\label{sec:orgd40e8f8}
El telescopio de Galileo, el teclado de Frescobaldi, la orquesta de
Monteverdi, el método inductivo de Bacon, el tratado de Descartes, la
fusión de estilos de Schütz: todos son tecnologías. Unas miran al
cielo; otras producen sonido; otras organizan el pensamiento. Pero
todas comparten una misma actitud: la materia no es un misterio
contemplativo, sino un campo de fuerzas que puede ser explorado,
medido y transformado.

Esta actitud —la actitud experimental— es la que, tres siglos después,
animará la ética hacker. El hacker no contempla el código: lo modifica.
El músico barroco no contemplaba la partitura: la improvisaba. El
científico del siglo XVII no contemplaba el cosmos: lo medía. En los
tres casos, la tecnología no es un medio pasivo: es una extensión del
cuerpo y del pensamiento.
\end{enumerate}
\subsubsection{El siglo XVIII: cálculo, contrapunto y armonía}
\label{sec:orgd46873c}
El siglo XVIII fue el siglo de la razón. La ciencia perfeccionó el
telescopio —con los reflectores de Newton y Herschel— y abrazó el
cálculo infinitesimal. La música abandonó la polifonía renacentista y
abrazó la armonía tonal. Dos herramientas —el cálculo y el bajo
continuo— transformaron la forma de entender el universo y el sonido.
\begin{enumerate}
\item Newton y Leibniz: el cálculo como lenguaje
\label{sec:orgaf4cc08}
Isaac Newton (1643-1727) y Gottfried Wilhelm Leibniz (1646-1716)
desarrollaron el cálculo infinitesimal de forma independiente. La
disputa por la prioridad —que amargó los últimos años de ambos— no
debe ocultar lo esencial: el cálculo es la primera herramienta
matemática que permite describir el cambio continuo. Una órbita
planetaria, una vibración sonora, una onda que se propaga: todo eso
es cálculo.

Newton aplicó el cálculo al cosmos. Sus \textbf{Principia Mathematica} (1687)
unificaron la física terrestre y la celeste bajo una sola ley: la
gravitación universal. El universo era un mecanismo determinista
gobernado por ecuaciones. Leibniz aplicó el cálculo a la metafísica.
Su concepto de la \textbf{mónada} —una sustancia simple que refleja el
universo entero desde su perspectiva— era una ontología del punto de
vista. La música, para Leibniz, era «un ejercicio oculto de aritmética
del alma que no sabe que está contando». El alma calcula sin saberlo.
El músico hace matemáticas sin saberlo.

La tecnología de Newton y Leibniz era el papel, la tinta y la
correspondencia. Pero el cálculo era una tecnología en sí misma: una
herramienta mental que permitía pensar el cambio. Sin cálculo,
no hay análisis de Fourier. Y sin Fourier, no hay comprensión
científica del timbre. En 1822, Joseph Fourier demostró que cualquier
onda periódica —incluida la vibración de una cuerda, de una columna
de aire, de un parche— puede descomponerse en una suma de
ondas sinusoidales simples. Esas ondas son los armónicos o parciales:
la materia prima del sonido. Lo que Pitágoras intuyó con el monocordio
—que la octava, la quinta y la cuarta eran proporciones— Fourier lo
demostró con ecuaciones: el timbre de un violín, de una flauta, de una
campana, no es una cualidad misteriosa, sino una receta de armónicos.
Sin Fourier, no hay análisis espectral. Sin análisis espectral, no hay
síntesis digital. Sin síntesis digital, no hay SuperCollider. La
cadena es larga, pero el eslabón inicial está aquí.
\item Johann Sebastian Bach (1685-1750): el cálculo del contrapunto
\label{sec:orgb0f3581}
Bach no conocía el cálculo de Newton ni de Leibniz. Pero su música es
cálculo aplicado al sonido. Una fuga de Bach es un sistema de
ecuaciones donde cada voz es una variable que depende de las demás. El
contrapunto no es inspiración: es organización material de líneas
melódicas que se entrelazan según reglas precisas.

Bach trabajaba con la tecnología de su tiempo: el órgano de tubos, el
clavecín, el papel pautado. Pero su verdadera tecnología era la
notación musical: un sistema de signos que permitía escribir,
transmitir y transformar el sonido. Bach no improvisaba sus fugas: las
escribía. Y al escribirlas, las convertía en objetos de conocimiento.

La \textbf{Ofrenda Musical} (1747) es el ejemplo supremo. Bach improvisó un
tema a tres voces ante el rey Federico II de Prusia. Luego, en casa,
escribió una colección de cánones y fugas basados en ese tema,
explorando todas las posibilidades contrapuntísticas. Algunos cánones
están escritos de forma enigmática: solo una línea melódica, con una
indicación críptica de cómo debe ser interpretada. El intérprete debe
resolver el enigma. Bach no escribía música: escribía algoritmos.

\begin{verbatim}
\version "2.24.0"

\score {
  \new PianoStaff <<
    \new Staff {
      \clef "treble"
      \key c \major
      \time 4/4
      c''16 g' e' g' c'' g' e' g' d'' g' f' g' e'' g' c'' g'' |
      b'16 g' d' g' b' g' d' g' c'' g' e' g' d'' g' f' g' |
      a'16 f' c' f' a' f' c' f' b' f' d' f' c'' f' e' f' |
      g'16 e' c' e' g' e' c' e' a' e' d' e' b' e' c'' e' |
    }
    \new Staff {
      \clef "bass"
      \key c \major
      \time 4/4
      c,,1 g, a, f, |
      g,1 e, f, g, |
    }
  >>
  \layout {
    indent = 0
    ragged-right = ##f
  }
}
\end{verbatim}

\begin{figure}[htbp]
\centering
\includegraphics[width=.9\linewidth]{/home/spock/Documents/symploke-sonora/img/bach_preludio_bwv846.png}
\caption{J.S. Bach — Preludio en Do mayor, BWV 846}
\end{figure}

\begin{verbatim}
// Bach algorítmico: fuga a dos voces en SuperCollider
// El contrapunto como cálculo: cada voz es una función que depende de la otra

s.boot;

(
// SynthDef para una voz de la fuga (sonido de órgano simple)
SynthDef(\organo, { |out=0, freq=440, amp=0.1, gate=1|
    var sig, env;
    env = EnvGen.kr(Env.asr(0.02, 1, 0.1), gate, doneAction: 2);
    sig = Mix.ar([
        SinOsc.ar(freq, 0, 0.6),
        SinOsc.ar(freq * 2, 0, 0.2),
        SinOsc.ar(freq * 3, 0, 0.1)
    ]);
    sig = sig * env * amp;
    Out.ar(out, sig!2);
}).add;
)

// Esperar a que compile
0.5.wait;

// Sujeto de la fuga (tema principal) - BWV 846 adaptado
~sujeto = Pbind(
    \instrument, \organo,
    \midinote, Pseq([60, 64, 67, 72, 71, 67, 69, 65, 67, 64, 62, 60, 67, 65, 64, 62, 60], 1),
    \dur, Pseq([0.25, 0.25, 0.25, 0.5, 0.25, 0.25, 0.25, 0.25, 0.5, 0.25, 0.25, 0.5, 0.25, 0.25, 0.25, 0.25, 0.5], 1),
    \amp, 0.12
);

// Contrasujeto (segunda voz) - entra después, responde al sujeto
~contrasujeto = Pbind(
    \instrument, \organo,
    \midinote, Pseq([48, 52, 55, 60, 59, 55, 57, 53, 55, 52, 50, 48, 55, 53, 52, 50, 48], 1),
    \dur, Pseq([0.25, 0.25, 0.25, 0.5, 0.25, 0.25, 0.25, 0.25, 0.5, 0.25, 0.25, 0.5, 0.25, 0.25, 0.25, 0.25, 0.5], 1),
    \amp, 0.08
);

// Ejecutar las dos voces: la segunda entra después (stretto)
~voz1 = ~sujeto.play;
2.wait;  // La segunda voz entra dos segundos después
~voz2 = ~contrasujeto.play;

// Detener
~voz1.stop;
~voz2.stop;
)
\end{verbatim}

\begin{verbatim}
// Fuga a seis voces en SuperCollider
// Inspirado en Bach - cada voz entra en stretto

s.boot;

(
// SynthDef mejorado con más cuerpo para seis voces
SynthDef(\organo6, { |out=0, freq=440, amp=0.05, gate=1, pan=0|
    var sig, env;
    env = EnvGen.kr(Env.asr(0.03, 1, 0.15), gate, doneAction: 2);
    sig = Mix.ar([
        SinOsc.ar(freq, 0, 0.5),
        SinOsc.ar(freq * 2, 0, 0.15),
        SinOsc.ar(freq * 3, 0, 0.08),
        SinOsc.ar(freq * 4, 0, 0.04)
    ]);
    sig = sig * env * amp;
    Out.ar(out, Pan2.ar(sig, pan));
}).add;
)

0.5.wait;

(
// Sujeto principal (voz 1 - soprano)
~voz1 = Pbind(
    \instrument, \organo6,
    \midinote, Pseq([72, 76, 79, 84, 83, 79, 81, 77, 79, 76, 74, 72], 1),
    \dur, Pseq([0.5, 0.5, 0.5, 1.0, 0.5, 0.5, 0.5, 0.5, 1.0, 0.5, 0.5, 1.0], 1),
    \amp, 0.06,
    \pan, 0
);

// Respuesta real (voz 2 - contralto)
~voz2 = Pbind(
    \instrument, \organo6,
    \midinote, Pseq([60, 64, 67, 72, 71, 67, 69, 65, 67, 64, 62, 60], 1),
    \dur, Pseq([0.5, 0.5, 0.5, 1.0, 0.5, 0.5, 0.5, 0.5, 1.0, 0.5, 0.5, 1.0], 1),
    \amp, 0.05,
    \pan, -0.3
);

// Voz 3 (tenor) - entra en la quinta
~voz3 = Pbind(
    \instrument, \organo6,
    \midinote, Pseq([67, 71, 74, 79, 78, 74, 76, 72, 74, 71, 69, 67], 1),
    \dur, Pseq([0.5, 0.5, 0.5, 1.0, 0.5, 0.5, 0.5, 0.5, 1.0, 0.5, 0.5, 1.0], 1),
    \amp, 0.05,
    \pan, 0.3
);

// Voz 4 (bajo) - entra en la octava inferior
~voz4 = Pbind(
    \instrument, \organo6,
    \midinote, Pseq([48, 52, 55, 60, 59, 55, 57, 53, 55, 52, 50, 48], 1),
    \dur, Pseq([0.5, 0.5, 0.5, 1.0, 0.5, 0.5, 0.5, 0.5, 1.0, 0.5, 0.5, 1.0], 1),
    \amp, 0.06,
    \pan, 0
);

// Voz 5 - contrasujeto libre
~voz5 = Pbind(
    \instrument, \organo6,
    \midinote, Pseq([64, 62, 60, 59, 57, 60, 62, 64, 65, 67, 69, 67], 1),
    \dur, Pseq([0.25, 0.25, 0.5, 0.5, 0.5, 0.5, 0.5, 0.5, 0.5, 0.5, 0.5, 1.0], 1),
    \amp, 0.04,
    \pan, -0.6
);

// Voz 6 - pedal grave
~voz6 = Pbind(
    \instrument, \organo6,
    \midinote, Pseq([36, 36, 36, 36, 43, 43, 43, 43, 48, 48, 48, 48], 1),
    \dur, Pseq([1.0, 1.0, 1.0, 1.0, 1.0, 1.0, 1.0, 1.0, 1.0, 1.0, 1.0, 1.0], 1),
    \amp, 0.07,
    \pan, 0.6
);

// Iniciar las seis voces en stretto
("=== Fuga a seis voces ===").postln;
("Voz 1: Soprano").postln;
~p1 = ~voz1.play;

1.5.wait;
("Voz 2: Contralto (respuesta)").postln;
~p2 = ~voz2.play;

1.5.wait;
("Voz 3: Tenor").postln;
~p3 = ~voz3.play;

1.5.wait;
("Voz 4: Bajo").postln;
~p4 = ~voz4.play;

1.5.wait;
("Voz 5: Contrasujeto libre").postln;
~p5 = ~voz5.play;

1.5.wait;
("Voz 6: Pedal grave").postln;
~p6 = ~voz6.play;
)

// Para detener todas las voces:
// ~p1.stop; ~p2.stop; ~p3.stop; ~p4.stop; ~p5.stop; ~p6.stop;
\end{verbatim}
\item Jean-Philippe Rameau (1683-1764): la armonía como ciencia
\label{sec:orga104e04}
Rameau fue el Newton de la música. Así como Newton unificó la física
terrestre y la celeste, Rameau unificó la práctica musical y la teoría
armónica. Su \textbf{Traité de l'harmonie} (1722) no es un manual de
composición: es un intento de reducir la música a principios
naturales.

Partió de una observación empírica: la serie de armónicos. Todo sonido
contiene en sí mismo una serie de frecuencias secundarias —los
armónicos— que determinan su timbre. Esa serie no es arbitraria: sigue
proporciones matemáticas fijas. La armonía no es una convención
cultural: es una propiedad física del sonido. Tres conceptos resumen
su revolución: el bajo fundamental, la inversión de los acordes y el
acorde como unidad básica de la música. El bajo fundamental es el
principio que Rameau descubrió en 1722: toda armonía puede reducirse a
una sucesión de acordes cuyas fundamentales se mueven por intervalos
de quinta o tercera.

Rameau fue criticado en su tiempo. Los defensores de la tradición
contrapuntística —los "Lullystes"— lo acusaron de ser un "mero
teórico", a pesar de que Rameau era también un gran compositor de
óperas. La lucha entre los Lullystes y los "Rameauneurs" fue una lucha
de paradigmas. Y como en toda lucha de paradigmas, la academia se
resistió al cambio. Rameau no vivió para ver la aceptación plena de su
sistema, pero su \textbf{Traité} se convirtió en la base de la enseñanza
armónica hasta nuestros días.

\begin{verbatim}
// Rameau: bajo fundamental y enlace de acordes
// Traité de l'harmonie, 1722 — SuperCollider, 2026

s.boot;

(
// SynthDef para cada nota del acorde
SynthDef(\rameau, { |out=0, freq=440, amp=0.08, gate=1|
    var sig, env;
    env = EnvGen.kr(Env.asr(0.02, 1, 0.3), gate, doneAction: 2);
    sig = Mix.ar([
        SinOsc.ar(freq, 0, 0.5),
        SinOsc.ar(freq * 2, 0, 0.2),
        SinOsc.ar(freq * 3, 0, 0.1)
    ]);
    sig = sig * env * amp;
    Out.ar(out, sig!2);
}).add;
)

(
// Progresión de acordes según el bajo fundamental
// I (Do mayor) -> V (Sol mayor) -> I (Do mayor)
// Fundamentales: Do -> Sol -> Do (movimiento de quinta)
~progresion_rameau = Routine({
    var acordes = [
        [60, 64, 67],  // Do mayor (C-E-G) - Tónica
        [67, 71, 74],  // Sol mayor (G-B-D) - Dominante
        [60, 64, 67],  // Do mayor (C-E-G) - Tónica
        [60, 65, 69],  // La menor (A-C-E) - Subdominante relativa
        [67, 71, 74],  // Sol mayor (G-B-D) - Dominante
        [60, 64, 67]   // Do mayor (C-E-G) - Tónica
    ];
    var duraciones = [2.0, 1.5, 2.0, 1.5, 1.5, 3.0];

    ("=== Rameau: bajo fundamental ===").postln;
    acordes.do { |notas, i|
        var nombre = switch(i,
            0, {"I — Tónica (Do mayor)"},
            1, {"V — Dominante (Sol mayor)"},
            2, {"I — Tónica (Do mayor)"},
            3, {"vi — Subdominante relativa (La menor)"},
            4, {"V — Dominante (Sol mayor)"},
            5, {"I — Tónica (Do mayor) — cierre"}
        );
        ("→ " ++ nombre).postln;

        // Tocar el acorde: tres notas simultáneas
        notas.do { |midi|
            Synth(\rameau, [\freq, midi.midicps, \amp, 0.06]);
        };

        duraciones[i].wait;
    };
    ("=== El bajo fundamental organiza la armonía ===").postln;
});

~progresion_rameau.play;
)
\end{verbatim}
\item La Tensión entre Fux y Rameau: Paradigma y Dialéctica
\label{sec:org188649d}
La publicación del \textbf{Traité de l'harmonie} de Rameau en 1722 y del
\textbf{Gradus ad Parnassum} de Fux en 1725 no fue una coincidencia: fue el
choque de dos paradigmas. Fux representa la culminación de la
tradición contrapuntística: la música como arte de combinar líneas
melódicas independientes. Su método, basado en el estilo de
Palestrina, fue utilizado por Bach, Haydn, Mozart y Beethoven como
fundamento de su oficio. Rameau representa una nueva forma de entender
la música: la armonía como principio fundamental, la homofonía como
textura donde una voz principal se apoya en un acompañamiento
acórdico.

Esta tensión entre polifonía y homofonía es un ejemplo claro de lo que
Thomas Kuhn llamó un \textbf{cambio de paradigma}. Durante siglos, el
paradigma dominante fue el contrapunto; Rameau propuso uno nuevo. Pero
este cambio no fue pacífico: la teoría de Rameau fue recibida con
hostilidad por los defensores de la tradición. La lucha entre los
"Lullystes" y los "Rameauneurs" fue una lucha de paradigmas.

Desde la perspectiva hegeliana, esta lucha no es un obstáculo; es el
motor de la historia. La tesis —la polifonía como paradigma
dominante— genera su antítesis —la homofonía como desafío— y de esa
contradicción surge una síntesis que no es la simple suma de ambas,
sino una transformación del campo musical. La música del Clasicismo y
del Romanticismo no es ni puramente polifónica ni puramente
homofónica; es una síntesis de ambas tensiones.

Así, la tensión entre Fux y Rameau no es un simple debate técnico; es
un ejemplo de cómo el conocimiento musical avanza a través de la
contradicción. La academia, al resistirse a las innovaciones de
Rameau, no hizo sino poner en marcha el proceso dialéctico que, con el
tiempo, llevaría a la aceptación de su sistema como fundamento de la
teoría armónica moderna. Y esa misma resistencia es la que hoy
encuentro en mi camino: la academia puede tardar en reconocer nuevos
saberes, pero la historia demuestra que el cambio, aunque lento, es
inevitable.
\item La tecnología que los unió: el piano forte
\label{sec:org7c170a1}
El siglo XVIII vio nacer una tecnología que transformaría la música
para siempre: el piano forte. A diferencia del clavecín —que pellizca
las cuerdas y no permite variar la intensidad—, el piano forte
golpea las cuerdas con martillos cuya fuerza depende de la pulsación
del dedo. Por primera vez, un teclado podía producir \textbf{piano} y \textbf{forte},
matices, gradaciones dinámicas.

El piano forte no era solo un instrumento: era una ontología. Permitía
pensar la música como un continuo de intensidades, no como una
sucesión de notas discretas. El \textbf{cromatofonismo} —el eje Z del volumen
tridimensional de Chuliá— encontraba en el piano forte su tecnología
más perfecta.

Bach probó los primeros pianos forte hacia el final de su vida. No le
gustaron: prefería el clavecín y el órgano. Rameau tampoco los usó
mucho. Pero Mozart, Haydn y Beethoven crecerían con el piano forte
como herramienta principal. La tecnología precede al pensamiento: el
piano forte hizo posible la música del Clasicismo y del Romanticismo.
\item Paralelismos y divergencias
\label{sec:orgf52846a}
Newton y Bach nunca se conocieron. Leibniz y Bach sí: Leibniz
frecuentaba la corte de Hannover, donde Bach trabajó brevemente. Pero
no hay registro de un encuentro significativo.

Sin embargo, la ciencia de Newton y Leibniz y la música de Bach y
Rameau comparten un mismo espíritu: la búsqueda de leyes universales
que expliquen la diversidad de los fenómenos. La gravitación universal
explica la caída de una manzana y la órbita de un planeta. La armonía
tonal explica un acorde de Do mayor y una sinfonía de Haydn. En ambos
casos, la ley es simple; las consecuencias, infinitas.

Este espíritu —la confianza en la razón, la búsqueda de principios
universales, la fe en el progreso— es el espíritu de la Ilustración.
Y la Ilustración es, también, el espíritu del hacker: el código es
simple; las aplicaciones, infinitas. El cálculo de Newton y Leibniz es
el antepasado del algoritmo. La fuga de Bach es el antepasado del
programa. La armonía de Rameau es el antepasado del sintetizador.
\end{enumerate}
\subsubsection{El Clasicismo vienés: Haydn, Mozart, Beethoven}
\label{sec:org310275f}
Si el Barroco fue el siglo del cálculo, el Clasicismo fue el siglo de
la forma. La música no solo debía ser bella: debía ser inteligible.
La forma sonata —exposición, desarrollo, reexposición— se convirtió en
la arquitectura musical por excelencia. Tres compositores, los tres
en Viena, llevaron esa arquitectura a su culminación.
\begin{enumerate}
\item Joseph Haydn (1732-1809): el inventor de la sinfonía
\label{sec:org3fb3420}
Haydn no inventó la sinfonía, pero la convirtió en lo que hoy
conocemos. Trabajó durante décadas para la familia Esterházy, aislado
en un palacio rural, con una orquesta a su disposición. Ese
aislamiento —esa libertad material— le permitió experimentar. «Tuve
que ser original», escribió. La tecnología de Haydn era la orquesta
clásica: cuerdas, maderas, metales, timbales. Un laboratorio de
timbres que él organizó como quien diseña un edificio.
\item Wolfgang Amadeus Mozart (1756-1791): la ópera como laboratorio
\label{sec:org4e8ba1d}
Mozart fue el primer compositor que intentó vivir como freelance, sin
mecenas. Fracasó. Murió en la pobreza. Pero en sus partituras no hay
rastro de ese fracaso: \textbf{Eine kleine Nachtmusik}, K. 525 (1787), es una
ecuación resuelta con sonidos. Simetría, equilibrio, perfección
formal. La tecnología de Mozart era el piano forte y la orquesta,
pero sobre todo era la ópera: un género que unificaba música, texto,
escenografía y público. En \textbf{Don Giovanni} (1787) y \textbf{La flauta mágica}
(1791), exploró los límites de la armonía tonal al servicio del drama.

Esta disociación entre la vida material y la obra es una lección
ontológica. Mozart no expresaba su sufrimiento en la música —eso lo
haría Beethoven después—: lo ocultaba bajo una superficie de
elegancia. La materia sonora (\(M_1\)) no refleja necesariamente la
materia biográfica (\(M_2\)). La música no es un diario íntimo: es una
organización de vibraciones que puede ser independiente de las
emociones de quien la produce. El idealismo romántico quiere que el
arte exprese el alma del artista. Mozart demuestra lo contrario.

\begin{verbatim}
\version "2.24.0"

\score {
  \new StaffGroup <<
    \new Staff {
      \clef "treble"
      \key g \major
      \time 4/4
      \partial 4
      d''8.\p d''16 |
      d''4 r8 d''8 e''8. d''16 d''4 |
      r8 d''8 g''8. fis''16 e''8. d''16 c''8. b'16 |
      a'4 r8 a'8 b'8. a'16 a'4 |
    }
    \new Staff {
      \clef "bass"
      \key g \major
      \time 4/4
      \partial 4
      r4 |
      g,1 | g, | d, | d, |
    }
  >>
  \layout {
    indent = 0
    ragged-right = ##f
  }
}
\end{verbatim}

\begin{figure}[htbp]
\centering
\includegraphics[width=.9\linewidth]{/home/spock/Documents/symploke-sonora/img/mozart_nachtmusik.png}
\caption{W.A. Mozart — Eine kleine Nachtmusik, K. 525 (1787)}
\end{figure}
\item Ludwig van Beethoven (1770-1827): la subjetividad como forma
\label{sec:orgad8823d}
Beethoven no rompió con el Clasicismo: lo llevó a su límite. Su
tecnología seguía siendo la orquesta y el piano forte. Pero su método
era nuevo: la forma sonata ya no era un molde, sino un campo de
batalla. En la \textbf{Quinta Sinfonía} (1808), cuatro notas —tres breves,
una larga— se convierten en el motor de toda la obra. La forma no es
un recipiente: es un organismo que crece.

Beethoven fue también el primer compositor que usó sistemáticamente el
metrónomo, patentado por Maelzel en 1815. Pero el propio Beethoven
desconfiaba de su invento. En una carta al director de orquesta Ignaz
von Mosel, fechada en Viena en diciembre de 1817, criticó el uso
rígido del metrónomo y defendió que las indicaciones tradicionales de
tempo —*Allegro*, *Andante*— describen el carácter de la obra, no una
velocidad mecánica. La música, incluso para el compositor que abrazó
el metrónomo, es más flexible que el determinismo de un artefacto.
\item Kant, Hegel y la música del Clasicismo
\label{sec:org0cc1e94}
El Clasicismo vienés coincide con el idealismo alemán. Kant publicó la
\textbf{Crítica del Juicio} en 1790, cuando Mozart componía \textbf{Così fan tutte}.
Hegel dictó sus \textbf{Lecciones de Estética} en la década de 1820, cuando
Beethoven escribía sus últimos cuartetos. Ambos pensaron la música
desde la filosofía, pero con resultados opuestos.

Kant, como vimos, consideró la música un "juego de emociones" inferior
a la poesía. No entendió la forma sonata. No entendió que la música
del Clasicismo no era un juego: era una arquitectura. Hegel, en
cambio, situó la música en un lugar ambiguo pero elevado: es la más
subjetiva de las artes, la que expresa el sentimiento interior, pero
también la más vacía de contenido cognitivo. La música, para Hegel, es
el espíritu que se experimenta a sí mismo sin mediación de la imagen
ni de la palabra.

Desde el materialismo filosófico, tanto Kant como Hegel comparten un
mismo error: juzgan la música desde el concepto (\(M_3\)), no desde la
materia (\(M_1\)). La forma sonata no es una idea: es una organización
material de vibraciones en el tiempo. Beethoven no expresaba
sentimientos: organizaba sonidos. Que esos sonidos produzcan
sentimientos en el oyente (\(M_2\)) es un efecto, no una causa.
\item La tecnología del Clasicismo: el piano forte y la orquesta
\label{sec:org1a31b3b}
El piano forte, que Bach había rechazado, se convirtió en el
instrumento central del Clasicismo. Su capacidad para producir matices
dinámicos —*piano* y \textbf{forte}, de ahí su nombre— permitió una
expresividad que el clavecín no podía ofrecer. La orquesta clásica,
por su parte, se estandarizó: cuerdas, maderas a dos, metales,
timbales. Una formación que permitía claridad, equilibrio, contraste.

Pero estas tecnologías no eran neutrales: condicionaban lo que se
podía pensar y, sobre todo, lo que se podía escuchar. Vicente Chuliá
ha señalado que el piano —y más concretamente el piano moderno de
cola— es el instrumento que materializa una imposición violenta sobre
la materia sonora: el temperamento igual. Los doce semitonos no son
una evolución natural de la música, sino una decisión histórica que
falsea los intervalos. Antes del temperamento igual, una quinta era
exactamente 3:2; una tercera mayor, 5:4. Cada tonalidad sonaba
distinta. El temperamento igual igualó artificialmente todos los
intervalos para que las doce tonalidades fueran equivalentes, pero lo
hizo a costa de la pureza armónica. Ninguna quinta en un piano
moderno es exactamente 3:2. El oído lo sabe, aunque la costumbre lo
silencie.

El piano impone esa convención a todo lo que lo rodea. Un violinista
puede ajustar la afinación de cada nota; un cantante puede entonar una
tercera justa; un pianista no. Cuando una orquesta afina con un piano,
se pliega a su temperamento. Por eso el oboe da el La: porque el oboe
puede ajustarse a la orquesta; el piano, no. Lo que se ganó en
versatilidad —modular libremente entre tonalidades— se perdió en
riqueza espectral.

La comparación con la Revolución Francesa, sugerida por Chuliá, es
precisa: el piano impuso los doce semitonos del mismo modo que la
Revolución impuso la igualdad jurídica. Ambos fueron proyectos
ilustrados: racionales, sistemáticos, niveladores. Ambos falsearon
diferencias reales en nombre de un principio abstracto. La lección
ontológica es clara: confundir el instrumento con la realidad —creer
que la música tiene doce notas porque el piano tiene doce teclas— es
el pecado original del idealismo. La música tiene infinitas alturas;
el piano solo doce. La materia (\(M_1\)) desborda siempre la convención
(\(M_3\)).
\end{enumerate}
\subsubsection{Ludwig van Beethoven (1770-1827): la subjetividad como forma}
\label{sec:org2da984a}
Beethoven no rompió con el Clasicismo: lo llevó a su límite. Su
tecnología seguía siendo la orquesta y el piano forte. Pero su método
era nuevo: la forma sonata ya no era un molde, sino un campo de
batalla. En la \textbf{Quinta Sinfonía} (1808), cuatro notas —tres breves,
una larga— se convierten en el motor de toda la obra. La forma no es
un recipiente: es un organismo que crece.

Beethoven fue también el primer compositor que usó sistemáticamente el
metrónomo, patentado por Maelzel en 1815. Pero el propio Beethoven
desconfiaba de su invento. En una carta al director de orquesta Ignaz
von Mosel, fechada en Viena en diciembre de 1817, criticó el uso
rígido del metrónomo y defendió que las indicaciones tradicionales de
tempo —*Allegro*, *Andante*— describen el carácter de la obra, no una
velocidad mecánica. La música, incluso para el compositor que abrazó
el metrónomo, es más flexible que el determinismo de un artefacto.

Beethoven no fue un romántico. Su formación era clásica: contrapunto,
forma sonata, armonía tonal. Nunca se llamó a sí mismo "romántico".
Pero su tercera etapa —los últimos cuartetos de cuerda, la \textbf{Sinfonía
nº 9} (1824), la \textbf{Missa Solemnis*— anticipó elementos que serían
centrales en el Romanticismo: la expresión subjetiva extrema, la
disolución de la forma clásica, la introducción de la voz humana en la
sinfonía, el uso del silencio como material expresivo. La *Grosse
Fuge} (1826) era tan radical que los editores le pidieron un final
alternativo. Stravinsky dijo de ella: «Es una obra absolutamente
contemporánea que será contemporánea para siempre».

Beethoven no rompió con la tradición: la llevó a un punto de tensión
tal que sus sucesores no pudieron continuarla sin transformarla. Es un
cierre categorial: la culminación del Clasicismo que, al alcanzar su
límite, abre el camino al Romanticismo.
\begin{enumerate}
\item El féretro y la antorcha: Schubert
\label{sec:org970f639}
El 29 de marzo de 1827, Beethoven fue enterrado en el cementerio de
Währing, en Viena. Entre los portadores del féretro había un joven de
treinta años, admirador devoto del maestro: Franz Schubert
(1797-1828). Schubert murió un año después, en 1828, y fue enterrado
junto a Beethoven, como había pedido.

Ese féretro es un cierre categorial en sí mismo: el Clasicismo
muere con Beethoven, pero el Romanticismo nace en el mismo gesto.
Schubert ya había compuesto más de seiscientas canciones —los
\textbf{Lieder*— que transformaban el poema en experiencia subjetiva. Su
*Sinfonía Inacabada} (1822) y su \textbf{Quinteto para cuerdas en Do mayor}
(1828) no son ejercicios de forma: son exploraciones del sentimiento,
del paisaje interior, de la melancolía. Schubert no escribía
arquitectura sonora: escribía diarios íntimos con sonido.

La imagen es poderosa: el clásico que muere, el romántico que carga su
féretro. Dos mundos que se tocan en un gesto. La \textbf{symploké} hecha
ritual funerario.
\end{enumerate}
\subsubsection{El Romanticismo: entre el sentimiento y la máquina}
\label{sec:orgbdb2dca}
El Romanticismo musical suele presentarse como una rebelión del
sentimiento contra la razón. Pero esa narrativa es engañosa. El siglo
XIX fue también el siglo de la máquina de vapor, del ferrocarril, del
telégrafo, de la termodinámica, del electromagnetismo. Los mismos
compositores que exaltaban la subjetividad usaban tecnologías cada vez
más sofisticadas: orquestas gigantescas, instrumentos de válvulas,
pianos de cola con estructuras de hierro. El Romanticismo no fue una
huida de la técnica: fue una negociación entre la emoción y la
máquina.
\begin{enumerate}
\item Hector Berlioz (1803-1869): la orquesta como laboratorio
\label{sec:orgde33fef}
Berlioz llevó la orquesta a dimensiones colosales. Su \textbf{Sinfonía
Fantástica} (1830) requería más de cien músicos, incluyendo cuatro
arpas, campanas y una sección de vientos expandida. No era
grandilocuencia: era experimentación tímbrica. Berlioz trataba la
orquesta como un sintetizador analógico: cada combinación de
instrumentos producía un color nuevo.

En 1844, Berlioz publicó el \textbf{Grand Traité d'Instrumentation et
d'Orchestration Modernes}, el primer tratado sistemático de
orquestación. No era un manual: era un catálogo de posibilidades
tímbricas. Cada instrumento era descrito no por su historia, sino por
sus capacidades materiales: qué notas podía tocar, qué dinámicas
producía, qué efectos especiales permitía. Berlioz no preguntaba "¿qué
significa este instrumento?", sino "¿qué puede hacer este
instrumento?". Esa pregunta es materialista.

\begin{verbatim}
library(ggplot2)

orquesta <- data.frame(
  seccion = c("Cuerdas", "Maderas", "Metales", "Percusión", "Arpas"),
  cantidad = c(60, 10, 13, 4, 4)
)

ggplot(orquesta, aes(x = reorder(seccion, -cantidad), y = cantidad, fill = seccion)) +
  geom_col(width = 0.6, color = "white", linewidth = 0.5) +
  geom_text(aes(label = paste0(cantidad, " músicos")), vjust = -0.5, size = 5, fontface = "bold") +
  scale_fill_manual(values = c(
    "Cuerdas" = "#922B21",
    "Maderas" = "#1A5276", 
    "Metales" = "#D4AC0D",
    "Percusión" = "#6C3483",
    "Arpas" = "#117864"
  )) +
  labs(title = "Sinfonía Fantástica — Hector Berlioz (1830)",
       subtitle = "Más de 90 músicos. La orquesta como laboratorio tímbrico.",
       x = "", y = "Número de músicos") +
  theme_minimal(base_size = 16) +
  theme(axis.text.x = element_text(face = "bold", size = 14),
        plot.title = element_text(face = "bold", size = 20, hjust = 0.5),
        plot.subtitle = element_text(color = "gray40", size = 14, hjust = 0.5),
        legend.position = "none")
\end{verbatim}

\begin{figure}[htbp]
\centering
\includegraphics[width=.9\linewidth]{/home/spock/Documents/symploke-sonora/img/berlioz_orquesta.png}
\caption{Distribución de la orquesta en la Sinfonía Fantástica de Berlioz}
\end{figure}
\item Franz Liszt (1811-1886): el virtuosismo como tecnología
\label{sec:org570004b}
Liszt fue el primer virtuoso moderno. Su técnica pianística —escalas
en octavas, saltos imposibles, arpegios que cubrían todo el teclado—
no era un don natural: era una tecnología corporal desarrollada con
horas de práctica. Liszt no tocaba el piano: lo hackeaba.

En 1847, Liszt se retiró de los escenarios y se dedicó a la
composición y a la dirección. Sus \textbf{Poemas sinfónicos} —obras
orquestales basadas en un argumento extramusical— rompieron con la
forma sonata. La música no debía seguir una arquitectura abstracta,
sino un programa narrativo. Era un intento de unificar música y
literatura, sonido y significado.
\item Richard Wagner (1813-1883): la obra de arte total
\label{sec:org5887ff9}
Wagner llevó el Romanticismo a su culminación y a su crisis. Su
concepto de \textbf{Gesamtkunstwerk} —obra de arte total— pretendía unificar
música, poesía, teatro, escenografía y arquitectura en una sola
experiencia. Para ello construyó un teatro en Bayreuth, financiado por
el rey Luis II de Baviera, con un foso de orquesta oculto y una
acústica diseñada para envolver al espectador.

La tecnología de Wagner era el teatro mismo: una máquina de inmersión
sensorial. En \textbf{Tristán e Isolda} (1865), llevó la armonía tonal hasta
un punto de tensión insostenible. El acorde de Tristán —fa, si, re\#,
sol\#— no resuelve: se desliza hacia otra tensión. La tonalidad se
disuelve. Es el preludio del atonalismo que Schoenberg sistematizará
medio siglo después.
\item Mijaíl Bakunin (1814-1876) y Wagner: la revolución y la música
\label{sec:org6c333b8}
Wagner no fue solo un compositor: fue un revolucionario. Participó en
el levantamiento de Dresde de 1849 junto al anarquista ruso Mijaíl
Bakunin. Mientras Bakunin defendía las barricadas, Wagner imprimía
panfletos y hacía granadas de mano. Bakunin fue arrestado y enviado a
Siberia; Wagner escapó a Suiza.

La influencia de Bakunin sobre Wagner fue profunda. El \textbf{Anillo del
nibelungo} —cuya versión definitiva termina con la destrucción del
Walhalla, el palacio de los dioses— está impregnado de anarquismo
bakuniniano: el poder corrompe, el orden es violencia, la destrucción
es liberación. Pero Wagner también fue un nacionalista alemán y un
antisemita virulento. Su obra fue apropiada por el nazismo. La música
no es neutral: puede ser liberación y puede ser opresión.
\item Ciencia y tecnología en el siglo XIX
\label{sec:orgc81db3d}
Mientras los románticos exaltaban el sentimiento, la ciencia avanzaba
a un ritmo sin precedentes:

\begin{itemize}
\item \textbf{\textbf{Michael Faraday (1791-1867):}} Descubrió la inducción
electromagnética. Demostró que la electricidad y el magnetismo son
dos manifestaciones del mismo fenómeno. Su tecnología era el
laboratorio y el cuaderno de notas.

\item \textbf{\textbf{James Clerk Maxwell (1831-1879):}} Unificó la electricidad, el
magnetismo y la luz en un solo sistema de ecuaciones. La luz, dijo
Maxwell, es una onda electromagnética. Y si la luz es una onda,
también lo es el sonido. La música y la física comparten ontología.

\item \textbf{\textbf{Heinrich Hertz (1857-1894):}} Demostró experimentalmente la
existencia de las ondas electromagnéticas predichas por Maxwell. Sin
Hertz, no hay radio, no hay televisión, no hay transmisión
inalámbrica. La música que hoy escuchamos por streaming es posible
gracias a Hertz.

\item \textbf{\textbf{Joseph Fourier (1768-1830):}} Había demostrado en 1822 que
cualquier onda periódica puede descomponerse en una suma de ondas
sinusoidales. Los armónicos no son una intuición pitagórica: son una
realidad matemática. El timbre de un violín, de una flauta, de una
campana, es una receta de armónicos.

\begin{verbatim}
library(ggplot2)
library(dplyr)
library(tidyr)

# Crear una onda cuadrada como suma de armónicos impares
t <- seq(0, 2*pi, length.out = 1000)
f0 <- 1  # Frecuencia fundamental

# Onda cuadrada: suma de senos de armónicos impares con amplitud 1/n
onda_cuadrada <- sin(t * f0)
for(n in seq(3, 15, by = 2)) {
  onda_cuadrada <- onda_cuadrada + (1/n) * sin(t * f0 * n)
}
onda_cuadrada <- onda_cuadrada * (4/pi)  # Normalizar

# Datos para graficar
datos_onda <- data.frame(
  tiempo = t,
  amplitud = onda_cuadrada,
  tipo = "Onda resultante\n(suma de armónicos impares)"
)

# Armónicos individuales
armonicos <- data.frame(
  tiempo = rep(t, 4),
  armonico = rep(c("Fundamental (1f)", "3er armónico (3f)", "5to armónico (5f)", "7mo armónico (7f)"), each = 1000),
  amplitud = c(
    (4/pi) * sin(t * f0),
    (4/pi) * (1/3) * sin(t * f0 * 3),
    (4/pi) * (1/5) * sin(t * f0 * 5),
    (4/pi) * (1/7) * sin(t * f0 * 7)
  ),
  tipo = "Armónicos individuales"
)

# Combinar datos
todos_datos <- bind_rows(datos_onda, armonicos)

# Paleta de colores
colores <- c(
  "Fundamental (1f)" = "#1A5276",
  "3er armónico (3f)" = "#922B21", 
  "5to armónico (5f)" = "#D4AC0D",
  "7mo armónico (7f)" = "#6C3483"
)

# Gráfico de armónicos individuales
p1 <- ggplot(armonicos, aes(x = tiempo, y = amplitud, color = armonico)) +
  geom_line(linewidth = 0.8) +
  scale_color_manual(values = colores) +
  facet_wrap(~armonico, ncol = 2, scales = "free_y") +
  labs(title = "Armónicos individuales de una onda cuadrada",
       subtitle = "Cada armónico es una onda sinusoidal pura — Fourier, 1822",
       x = "Tiempo", y = "Amplitud") +
  theme_minimal(base_size = 12) +
  theme(legend.position = "none",
        strip.text = element_text(face = "bold", size = 10),
        plot.title = element_text(face = "bold", size = 16),
        plot.subtitle = element_text(color = "gray40"))

# Gráfico de la onda resultante
p2 <- ggplot(datos_onda, aes(x = tiempo, y = amplitud)) +
  geom_line(color = "#1A5276", linewidth = 1.2) +
  labs(title = "Onda cuadrada: suma de armónicos impares",
       subtitle = expression(frac(4, pi) * (sin(x) + frac(1, 3) * sin(3*x) + frac(1, 5) * sin(5*x) + ...)),
       x = "Tiempo", y = "Amplitud") +
  theme_minimal(base_size = 12) +
  theme(plot.title = element_text(face = "bold", size = 16),
        plot.subtitle = element_text(color = "gray40"))

# Mostrar ambos gráficos
library(patchwork)
p2 / p1 +
  plot_annotation(
    title = "Descomposición de Fourier: la base matemática del timbre",
    subtitle = "Joseph Fourier demostró que cualquier onda periódica puede descomponerse en una suma de ondas sinusoidales",
    theme = theme(plot.title = element_text(face = "bold", size = 18, hjust = 0.5),
                  plot.subtitle = element_text(color = "gray40", size = 12, hjust = 0.5))
  )
\end{verbatim}
\end{itemize}

\begin{figure}[htbp]
\centering
\includegraphics[width=.9\linewidth]{/home/spock/Documents/symploke-sonora/img/fourier_armonicos.png}
\caption{Descomposición de Fourier — Cualquier onda periódica es una suma de armónicos}
\end{figure}

\begin{verbatim}
// Descomposición de Fourier en SuperCollider
// Una onda cuadrada se construye sumando armónicos impares
// Fourier, 1822 — SuperCollider, 2026

(
// Onda cuadrada por síntesis aditiva
SynthDef(\cuadrada_fourier, { |out=0, freq=220, n_arm=1, amp=0.15|
    var sig;
    // Suma de senos: solo armónicos impares con amplitud 1/n
    sig = Mix.ar(
        Array.series(n_arm, 1, 2).collect { |n, i|
            SinOsc.ar(freq * n, 0, 1/n)
        }
    );
    sig = sig * (4/pi) * amp;  // Normalizar como Fourier
    Out.ar(out, sig!2);
}).add;

// Rutina para demostrar la construcción progresiva de la onda
~fourier_demo = Routine({
    var notas = [1, 2, 3, 5, 7, 10, 15, 20, 30, 50];
    ("=== Demostración de Fourier: onda cuadrada por suma de armónicos ===").postln;
    ("Frecuencia fundamental: 220 Hz (La3)").postln;
    notas.do { |n|
        var nombre;
        if(n == 1) { nombre = "Solo fundamental: seno puro" };
        if(n == 2) { nombre = "Fundamental + 3er armónico" };
        if(n > 2) { nombre = "Suma de " ++ n ++ " armónicos impares" };
        ("→ " ++ nombre).postln;
        Synth(\cuadrada_fourier, [\freq, 220, \n_arm, n]);
        2.wait;
    };
    ("=== La onda se aproxima a una cuadrada ===").postln;
});

// Ejecutar
~fourier_demo.play;
)

// Para detener: Cmd + . (o Ctrl + .)
\end{verbatim}

El mismo experimento puede realizarse en SuperCollider, activando
cada armónico individualmente.

\begin{verbatim}
// Laboratorio de Fourier: control de armónicos individuales
// Cada armónico es un parcial que puede activarse/desactivarse
// El timbre resultante es la suma de todos los parciales activos

s.boot;

(
SynthDef(\fourier_lab, { |out=0, freq=220, amp=0.08|
    var sig;
    sig = Mix.ar(16.collect { |i|
        var n = i + 1;
        var active = NamedControl.kr(("arm_" ++ n).asSymbol, 1, 0.02);
        SinOsc.ar(freq * n, 0, 1/n) * active
    });
    sig = sig * amp;
    Out.ar(out, sig!2);
}).add;
)

// Probar con fundamental sola (seno puro)
x = Synth(\fourier_lab, [\freq, 220, \amp, 0.15]);

// Activar armónicos uno a uno
x.set(\arm_2, 1);  // 2do armónico (octava)
x.set(\arm_3, 1);  // 3er armónico (quinta)
x.set(\arm_4, 1);  // 4to armónico (doble octava)
x.set(\arm_5, 1);  // 5to armónico (tercera mayor)

// Apagar la fundamental para escuchar solo armónicos
x.set(\arm_1, 0);

// Volver a encenderla
x.set(\arm_1, 1);

// Activar impares para aproximar onda cuadrada
x.set(\arm_7, 1);
x.set(\arm_9, 1);
x.set(\arm_11, 1);

// Liberar
x.free;
\end{verbatim}


\begin{itemize}
\item \textbf{\textbf{Urbain Le Verrier (1811-1877):}} Descubrió Neptuno en 1846 sin
mirar al cielo. Calculó su posición a partir de las perturbaciones
en la órbita de Urano. La astronomía ya no necesitaba telescopios:
necesitaba matemáticas. El cosmos se había convertido en un sistema
de ecuaciones.
\end{itemize}
\item Friedrich Nietzsche (1844-1900): la música como afirmación
\label{sec:org231e358}
Nietzsche fue músico antes que filósofo. Compuso canciones, piezas
para piano, intentos de ópera. Su admiración por Wagner fue devota
—llegó a llamarlo «mi maestro»—, pero terminó en ruptura. En \textbf{El caso
Wagner} (1888), acusó a Wagner de ser un «músico de la mentira» que
usaba el arte para manipular las emociones del público.

Para Nietzsche, la música no debía ser un consuelo ni una evasión,
sino una afirmación de la vida. Su ideal era la música dionisíaca: la
que celebra el devenir, el cuerpo, la embriaguez. Bizet —con su
*Carmen*— le parecía más saludable que Wagner. La música no debe
adoctrinar: debe liberar.

Pero Nietzsche no es un modelo para este ensayo. Fue un filólogo que
rompió con la academia, que renunció a su cátedra en Basilea a los
treinta y cuatro años, que escribió contra los sistemas filosóficos y
contra la universidad como institución. No buscó ser reconocido por
sus pares: buscó ser leído por un futuro que aún no existía. Su
trayectoria es la de un solitario que convirtió su vida en un
experimento y pagó el precio con la locura.

Mi camino es distinto. No busco romper con la academia: busco
contribuir a ella desde una tradición filosófica —el materialismo de
la Escuela de Toledo— que puede enriquecer sus líneas de
investigación. No aspiro al aislamiento del genio incomprendido, sino
a la discusión rigurosa entre pares. No escribo aforismos: propongo un
sistema abierto, enraizado en la experiencia y dispuesto a ser
confrontado. Nietzsche es un aliado en la crítica al idealismo, pero
no es un modelo de trayectoria. Mi apuesta es otra: seguir
aprendiendo, seguir investigando, seguir haciendo música-filosofía
mientras la academia me lo permita.
\item Paralelismos y divergencias
\label{sec:orgd74c096}
El Romanticismo musical y la ciencia del siglo XIX comparten un mismo
horizonte: la exploración de lo invisible. Maxwell demostró que la luz
es una onda; Berlioz demostró que la orquesta es un espectro. Hertz
detectó ondas electromagnéticas; Wagner disolvió la tonalidad. Fourier
descompuso el sonido en armónicos; Liszt descompuso el piano en
posibilidades.

Pero hay una diferencia fundamental: la ciencia del siglo XIX era
materialista; la música del Romanticismo era idealista. Los científicos
buscaban leyes; los músicos buscaban emociones. Los científicos
medían; los músicos sentían. Esta bifurcación —que Baumgarten había
inaugurado al separar la verdad de la belleza— alcanza en el siglo XIX
su punto máximo. La música ya no es una ciencia: es un arte. Y el arte
romántico, con su culto al genio y al sentimiento, se aleja cada vez
más de la materia.
\end{enumerate}
\subsubsection{El siglo XX: entre el control total y la indeterminación radical}
\label{sec:org45be74e}
El siglo XX fue el siglo de los extremos. La música se debatió entre
el control total —Schoenberg, el dodecafonismo, el serialismo
integral— y la indeterminación radical —Cage, el azar, el silencio—.
La ciencia, mientras tanto, desmontó las certezas del siglo XIX: la
mecánica cuántica reveló un mundo de probabilidades, la relatividad
curvó el espacio-tiempo, la teoría de la información digitalizó el
mundo. La música y la ciencia compartieron un mismo destino: la
pérdida de la certidumbre.
\begin{enumerate}
\item Arnold Schoenberg (1874-1951) y la crisis de la tonalidad
\label{sec:org915b115}
Schoenberg no fue un destructor: fue un sistematizador. Diagnosticó
que la tonalidad había llegado a su límite —Wagner y Debussy la habían
llevado hasta un punto de tensión insostenible— y propuso un nuevo
sistema: el dodecafonismo. Las doce notas de la escala cromática,
organizadas en una serie fija, sin jerarquías. Todas iguales. Ninguna
subordinada a otra.

El dodecafonismo era un sistema de control total sobre la materia
sonora. Pero a diferencia de la tonalidad —que se fundaba en la física
del sonido (\(M_1\)) y en la percepción logarítmica del oído (\(M_2\))—,
el dodecafonismo era puramente conceptual (\(M_3\)). Schoenberg intentó
imponer orden a una materia que se resistía a ser organizada de ese
modo. El resultado fue una música que rompió con la tradición pero que
no logró constituir un nuevo lenguaje compartido.

Desde el materialismo filosófico, el fracaso relativo del
dodecafonismo no es un argumento contra Schoenberg, sino una
confirmación de la discontinuidad de la materia. No basta con que un
sistema sea lógicamente impecable en \(M_3\): necesita entrelazarse con
\(M_1\) y \(M_2\). Schoenberg lo intentó. No lo logró. Pero su intento
abrió un campo de posibilidades que compositores como Webern, Boulez
y el propio Babbitt explorarían durante décadas.

\begin{verbatim}
// Schoenberg: serie dodecafónica sonificada
// Pierrot Lunaire, Op. 21 (1912) — SuperCollider, 2026

s.boot;

(
// SynthDef para cada nota de la serie
SynthDef(\serie, { |out=0, freq=440, amp=0.1, gate=1|
    var sig, env;
    env = EnvGen.kr(Env.perc(0.01, 0.8, 1, -4), gate, doneAction: 2);
    sig = SinOsc.ar(freq, 0, amp) * env;
    Out.ar(out, sig!2);
}).add;
)

(
// Rutina: tocar la serie en sus formas original, inversa y retrógrada
~serie_schoenberg = Routine({
    // Serie original P0: Do, Do#, Re, Re#, Mi, Fa, Fa#, Sol, Sol#, La, La#, Si
    var p0 = [60, 61, 62, 63, 64, 65, 66, 67, 68, 69, 70, 71];
    // Inversa: espejo alrededor de la primera nota
    var inversa = p0.collect { |n| 60 - (n - 60) };
    // Retrógrada: la serie al revés
    var retrograda = p0.reverse;

    ("=== Schoenberg: serie dodecafónica ===").postln;
    ("Serie original P0").postln;
    p0.do { |n| Synth(\serie, [\freq, n.midicps, \amp, 0.08]); 0.3.wait };
    1.wait;

    ("Serie inversa").postln;
    inversa.do { |n| Synth(\serie, [\freq, n.midicps, \amp, 0.06]); 0.3.wait };
    1.wait;

    ("Serie retrógrada").postln;
    retrograda.do { |n| Synth(\serie, [\freq, n.midicps, \amp, 0.06]); 0.3.wait };
    1.wait;

    ("Las doce notas sin jerarquía. Ninguna subordinada a otra.").postln;
});

~serie_schoenberg.play;
)
\end{verbatim}
\item Igor Stravinsky (1882-1971): el camaleón
\label{sec:orgf3620d2}
Stravinsky no siguió un camino lineal: cambió de estilo como quien
cambia de herramienta. El 29 de mayo de 1913, en el Théâtre des
Champs-Élysées de París, el público enloqueció. No era para menos: su
\textbf{Consagración de la primavera} superponía un acorde de Mi bemol mayor
con un Fa bemol mayor —politonalidad pura— y sus ritmos violentos
ningún oído educado en el Romanticismo podía procesar. Los bailarines
no oían la orquesta; Nijinsky gritaba los compases desde el lateral.
Stravinsky escapó por una ventana trasera. La música no volvería a ser
la misma.

Luego abrazó el neoclasicismo. Luego, ya en los años cincuenta, adoptó
el serialismo que había rechazado durante décadas. Stravinsky es la
prueba de que la música del siglo XX no fue una línea recta hacia la
atonalidad, sino un campo de fuerzas donde convivían tendencias
opuestas. Su famosa frase —«la tradición no es el culto a las
cenizas, sino la transmisión del fuego»— es un manifiesto
materialista: la tradición no es un museo, sino un laboratorio.

\begin{verbatim}
\version "2.24.0"

\score {
  \new PianoStaff <<
    \new Staff {
      \clef "treble"
      \time 2/4
      <ees' g' bes' e''>8 <ees' g' bes' e''> <ees' g' bes' e''> <ees' g' bes' e''> |
      <fes' aes' ces'' f''>8 <fes' aes' ces'' f''> <fes' aes' ces'' f''> <fes' aes' ces'' f''> |
    }
    \new Staff {
      \clef "bass"
      \time 2/4
      <ees, ees'>8 <ees, ees'> <ees, ees'> <ees, ees'> |
      <fes, fes'>8 <fes, fes'> <fes, fes'> <fes, fes'> |
    }
  >>
  \layout {
    indent = 0
    ragged-right = ##f
  }
}
\end{verbatim}

\begin{figure}[htbp]
\centering
\includegraphics[width=.9\linewidth]{/home/spock/Documents/symploke-sonora/img/stravinsky_consagracion.png}
\caption{I. Stravinsky — Acorde politonal de La consagración de la primavera (1913)}
\end{figure}

\begin{verbatim}
\version "2.24.0"

\score {
  \new PianoStaff <<
    \new Staff {
      \clef "treble"
      \time 2/4
      <ees' g' bes' e''>8 <ees' g' bes' e''> <ees' g' bes' e''> <ees' g' bes' e''> |
      <fes' aes' ces'' f''>8 <fes' aes' ces'' f''> <fes' aes' ces'' f''> <fes' aes' ces'' f''> |
    }
    \new Staff {
      \clef "bass"
      \time 2/4
      <ees, ees'>8 <ees, ees'> <ees, ees'> <ees, ees'> |
      <fes, fes'>8 <fes, fes'> <fes, fes'> <fes, fes'> |
    }
  >>
  \layout {
    indent = 0
    ragged-right = ##f
  }
}
\end{verbatim}
\item La mecánica cuántica y la crisis de la certeza
\label{sec:orge588927}
Mientras Schoenberg sistematizaba el dodecafonismo, la física
desmontaba el determinismo newtoniano. En 1927, Werner Heisenberg
(1901-1976) formuló el principio de incertidumbre: no se puede conocer
simultáneamente la posición y el momento de una partícula. No es una
limitación del instrumento de medición: es una propiedad de la
materia.

Ese mismo año, Niels Bohr (1885-1962) presentó el principio de
complementariedad: la luz es onda y partícula, dependiendo de cómo se
observe. La realidad no es una cosa: es una relación entre el
observador y lo observado.

Einstein, que había ayudado a fundar la cuántica, se resistió hasta el
final: «Dios no juega a los dados». Pero los dados estaban echados. La
materia, a nivel fundamental, es probabilidad, no certeza.
\item John Cage (1912-1992): la indeterminación como método
\label{sec:orga6b747d}
Cage estudió con Schoenberg. Pero llevó la lógica del control total
hasta su inversión: si Schoenberg quería controlarlo todo, Cage quiso
no controlar nada. Su obra \textbf{4'33''} (1952) —tres movimientos de
silencio— no es una broma: es la traducción musical del principio de
incertidumbre de Heisenberg. El compositor renuncia al control sobre
lo que suena; la música es el ambiente, el ruido, la respiración del
público.

Cage declaró que «no hay silencio». Siempre hay algo que suena. El
silencio no es ausencia: es presencia de lo impredecible. Su música no
es un mensaje: es un campo de posibilidades. Y esa es, precisamente,
la ontología de la mecánica cuántica.
\item La tecnología del siglo XX: del theremin al ordenador
\label{sec:org7702a94}
El siglo XX fue el siglo de la electrónica. En 1920, León Theremin
inventó el instrumento que lleva su nombre: dos antenas que detectan
la posición de las manos y convierten el movimiento en sonido. No hay
contacto físico: el cuerpo es el instrumento.

En los años cincuenta, Karlheinz Stockhausen (1928-2007) compuso
\textbf{Gesang der Jünglinge} (1956), combinando la voz de un niño con
sonidos electrónicos generados en el estudio de la WDR en Colonia. La
música electrónica no era un efecto especial: era un nuevo género.

En los años sesenta, Max Mathews (1926-2011) escribió MUSIC I, el
primer programa de síntesis digital. El ordenador se convirtió en
instrumento. En los años ochenta, Miller Puckette creó Pure Data. En
los noventa, James McCartney creó SuperCollider. La cadena es larga,
pero el eslabón inicial está aquí: el theremin, el estudio
electroacústico, el ordenador.

\begin{verbatim}
// John Cage: 4'33'' (1952)
// Tres movimientos de silencio. La música es lo que suena.

s.boot;

(
~cage_433 = Routine({
    ("=== John Cage: 4'33'' (1952) ===").postln;
    ("Primer movimiento: 33 segundos").postln;
    ("El compositor no controla el sonido.").postln;
    ("Escuche el ambiente. Eso es la música.").postln;
    33.wait;

    ("Segundo movimiento: 2 minutos 40 segundos").postln;
    ("No hay silencio. Siempre hay algo que suena.").postln;
    160.wait;

    ("Tercer movimiento: 1 minuto 20 segundos").postln;
    ("La música no es un objeto: es un proceso.").postln;
    80.wait;

    ("Fin de 4'33''. ¿Escuchó algo?").postln;
});

~cage_433.play;
)
\end{verbatim}
\item György Ligeti (1923-2006): la micropolifonía
\label{sec:org64d8e71}
Ligeti no siguió el camino de Schoenberg ni el de Cage. Su
micropolifonía —capas de voces tan densas que el oído no puede
distinguirlas— creaba texturas en movimiento, nubes sonoras que se
desplazaban lentamente. En \textbf{Atmosphères} (1961) y en el \textbf{Réquiem}
(1965), Ligeti demostró que la discontinuidad no es caos: es condición
del orden.

Su \textbf{Poème Symphonique} (1962) para cien metrónomos es la demostración
más elocuente de esta tesis. Los metrónomos se activan
simultáneamente, cada uno a un tempo diferente, y se detienen cuando
su cuerda se agota. Del caos inicial emerge un orden rítmico que
ninguno de los metrónomos contenía por sí solo. La materia, al
agotarse, revela sus estructuras subyacentes. No hay director, no hay
partitura: hay materia que se organiza.

\begin{verbatim}
library(ggplot2)

set.seed(1962)

n <- 40
tempos <- runif(n, 40, 208)
duraciones <- runif(n, 30, 180)

datos <- data.frame()
for(i in 1:n) {
  n_ticks <- floor(duraciones[i] / (60 / tempos[i]))
  if(n_ticks > 0) {
    ticks <- seq(0, duraciones[i], length.out = n_ticks)
    datos <- rbind(datos, data.frame(
      metronomo = i,
      tiempo = ticks,
      tempo = tempos[i]
    ))
  }
}

ggplot(datos, aes(x = tiempo, y = metronomo, color = tempo)) +
  geom_point(shape = "|", size = 0.5, alpha = 0.6) +
  scale_color_gradientn(colors = c("#1A5276", "#6C3483", "#922B21", "#D4AC0D"),
                        name = "Tempo (BPM)") +
  labs(title = "Poème Symphonique — György Ligeti (1962)",
       subtitle = "40 metrónomos. Sin director. Sin partitura. La materia se sincroniza al agotarse.",
       x = "Tiempo (segundos)", y = "Metrónomo nº") +
  theme(
    plot.background = element_rect(fill = "#1a1a2e", color = NA),
    panel.background = element_rect(fill = "#1a1a2e", color = NA),
    panel.grid.major = element_line(color = "gray20", linewidth = 0.2),
    panel.grid.minor = element_blank(),
    plot.title = element_text(face = "bold", size = 17, hjust = 0.5, color = "#eee"),
    plot.subtitle = element_text(color = "gray60", size = 11, hjust = 0.5),
    axis.title = element_text(color = "#ccc", size = 12),
    axis.text = element_text(color = "#aaa", size = 10),
    axis.text.y = element_blank(),
    axis.ticks.y = element_blank(),
    legend.background = element_rect(fill = "#1a1a2e", color = NA),
    legend.text = element_text(color = "#ccc"),
    legend.title = element_text(color = "#ccc"),
    legend.position = "bottom"
  )
\end{verbatim}

\begin{figure}[htbp]
\centering
\includegraphics[width=.9\linewidth]{/home/spock/Documents/symploke-sonora/img/ligeti_metronomos.png}
\caption{Simulación del Poème Symphonique de Ligeti}
\end{figure}
\item Iannis Xenakis (1922-2001): el arquitecto del sonido
\label{sec:orgb1da2b7}
Xenakis fue ingeniero civil antes que compositor. Colaboró con Le
Corbusier en el Pabellón Philips de la Exposición Universal de
Bruselas de 1958. Su música no era melodía ni armonía: era masa
sonora en movimiento. Su concepto de «música estocástica» aplicaba la
teoría de probabilidades a la composición: el compositor no controla
cada nota, sino la distribución estadística de los eventos.

En \textbf{Metastaseis} (1954), los glissandi de las cuerdas son superficies
que se deslizan, planos que se intersectan. La partitura de Xenakis no
es una sucesión de notas: es un plano arquitectónico. El eje
horizontal es el tiempo, el eje vertical es la altura. Xenakis no
componía para un espacio: componía con el espacio.

\begin{verbatim}
library(ggplot2)
library(viridis)

# Datos del paraboloide hiperbólico
x <- seq(-3, 3, length.out = 80)
y <- seq(-3, 3, length.out = 80)
grid <- expand.grid(x = x, y = y)
grid$z <- grid$x^2 - grid$y^2

ggplot(grid, aes(x = x, y = y, fill = z)) +
  geom_tile() +
  scale_fill_viridis(option = "magma", name = "Z = X² - Y²") +
  geom_contour(aes(z = z), color = "white", alpha = 0.3, linewidth = 0.3) +
  labs(title = "Paraboloide Hiperbólico — Pabellón Philips (1958)",
       subtitle = "La misma ecuación que Xenakis usó para Metastaseis (1954)",
       x = "X", y = "Y") +
  coord_fixed() +
  theme_minimal(base_size = 14) +
  theme(plot.background = element_rect(fill = "#1a1a2e", color = NA),
        panel.background = element_rect(fill = "#1a1a2e", color = NA),
        plot.title = element_text(face = "bold", size = 18, hjust = 0.5, color = "#eee"),
        plot.subtitle = element_text(color = "gray60", size = 12, hjust = 0.5),
        axis.text = element_text(color = "#aaa"),
        axis.title = element_text(color = "#ccc"),
        legend.background = element_rect(fill = "#1a1a2e", color = NA),
        legend.text = element_text(color = "#ccc"),
        legend.title = element_text(color = "#ccc"),
        panel.grid.major = element_line(color = "gray20", linewidth = 0.2),
        panel.grid.minor = element_blank())
\end{verbatim}

\begin{figure}[htbp]
\centering
\includegraphics[width=.9\linewidth]{/home/spock/Documents/symploke-sonora/img/xenakis_paraboloide.png}
\caption{Paraboloide hiperbólico — La geometría del Pabellón Philips}
\end{figure}
\begin{verbatim}
// Música estocástica inspirada en Xenakis
// Control estadístico: el compositor define la distribución,
// el algoritmo genera los eventos individuales
// Xenakis, Metastaseis, 1954

s.boot;

(
// SynthDef: cada "grano" es un evento sonoro independiente
SynthDef(\grano_xenakis, { |out=0, freq=440, amp=0.1, pan=0, duracion=0.1|
    var sig, env;
    env = EnvGen.kr(Env.perc(0.01, duracion, 1, -4), doneAction: 2);
    sig = SinOsc.ar(freq, 0, amp) * env;
    Out.ar(out, Pan2.ar(sig, pan));
}).add;
)

(
// Rutina estocástica: nube de sonidos controlada estadísticamente
~nube_estocastica = Routine({
    var densidad = 20;        // Eventos por segundo (textura)
    var centro_freq = 440;    // Frecuencia central (La4)
    var dispersion = 200;     // Dispersión alrededor del centro
    var duracion_base = 0.15; // Duración base de cada grano

    loop {
        // Frecuencia: distribución gaussiana alrededor del centro
        var freq = centro_freq + (dispersion * rrand(-1.0, 1.0)).gauss(0.5);
        // Amplitud: valores bajos para nubes suaves
        var amp = exprand(0.02, 0.08);
        // Panorámica: distribución espacial aleatoria
        var pan = rrand(-0.8, 0.8);
        // Duración variable de cada grano
        var dur = duracion_base * rrand(0.5, 2.0);

        Synth(\grano_xenakis, [
            \freq, freq.max(50).min(2000),
            \amp, amp,
            \pan, pan,
            \duracion, dur
        ]);

        (1.0 / densidad).wait;  // Control de densidad
    };
});

// Iniciar la nube estocástica
~nube_estocastica.play;
)

// Para detener: Cmd + . (o Ctrl + .)
\end{verbatim}
\item Milton Babbitt (1916-2011): el tiempo discontinuo
\label{sec:org0facdaa}
Babbitt fue matemático antes que compositor. Su \textbf{time-point system}
trata el tiempo musical no como un flujo continuo —un compás, un
pulso—, sino como una serie de puntos discretos, como píxeles en una
pantalla. El ritmo no es una sucesión de duraciones: es una
distribución de posiciones en una cuadrícula temporal.

En \textbf{Philomel} (1964), Babbitt llevó la síntesis electrónica más allá
de la mera experimentación tímbrica: la voz de la soprano es grabada,
analizada, fragmentada y recombinada electrónicamente. Toma una voz
humana (\(M_1\)), la convierte en datos (\(M_3\)) y la devuelve
transformada en sonido electrónico (\(M_1\) de nuevo). La fórmula de
Chuliá \(M_i \subset E \subset M\) encuentra en \textbf{Philomel} una de sus
realizaciones más perfectas.
\item Paralelismos y divergencias
\label{sec:org4defe2c}
El siglo XX fue el siglo de la bifurcación. La música académica se
alejó del público; el público se alejó de la música académica. La
ciencia se volvió incomprensible para el ciudadano medio; el ciudadano
medio se refugió en la tecnología sin entenderla. La especialización
—que había comenzado con Baumgarten— alcanzó su punto máximo.

Pero también fue el siglo de la convergencia. La electrónica unificó
la música y la física. El ordenador unificó el cálculo y la
creatividad. La información se convirtió en el nuevo paradigma. Y en
ese paradigma, el hacker —el músico-programador, el filósofo del
código— encontró su lugar.
\end{enumerate}
\subsubsection{El siglo XXI: el hacker como músico-filósofo}
\label{sec:orgde8516f}
El siglo XXI no empezó en el año 2000. Empezó en 1971, cuando Richard
Stallman —entonces un joven programador del MIT— se topó con una
impresora Xerox que se atascaba constantemente. Stallman quiso
modificar el código para que la impresora avisara a los usuarios
cuando se atascaba. Xerox se negó a darle el código fuente.

Ese gesto —negar el acceso al código— es el pecado original de la era
digital. Stallman respondió fundando la Free Software Foundation y
escribiendo la GNU General Public License: un documento legal que
garantiza cuatro libertades fundamentales: usar, estudiar, redistribuir
y mejorar el software. El software libre no es una cuestión técnica:
es una cuestión ética. Y la ética hacker —acceso libre, desconfianza
de la autoridad, mérito basado en la calidad del trabajo— es la
aplicación de esa ética al mundo digital.
\begin{enumerate}
\item El ordenador como instrumento musical
\label{sec:org59f9119}
El ordenador no fue diseñado para hacer música. Fue diseñado para
calcular trayectorias balísticas, descifrar códigos enemigos, procesar
datos. Pero como toda tecnología, una vez liberada de su propósito
original, encontró usos imprevistos.

En 1957, Max Mathews escribió MUSIC I en un IBM 704: el primer
programa de síntesis digital. El ordenador tardaba horas en generar
unos segundos de sonido. Pero el principio estaba establecido: el
cálculo podía convertirse en vibración. En los años ochenta, Miller
Puckette creó Pure Data, un entorno de programación gráfica para
síntesis de audio. En los años noventa, James McCartney creó
SuperCollider, un lenguaje de programación diseñado específicamente
para la síntesis y el procesamiento de audio en tiempo real.

SuperCollider no es un programa: es un laboratorio ontológico. Cada
línea de código es un experimento sobre \(M_1\) (la vibración), mediado
por \(M_2\) (la percepción en tiempo real) y organizado por \(M_3\) (el
concepto formalizado en algoritmos). Cuando escribo un sintetizador en
SuperCollider, estoy haciendo lo mismo que Rameau cuando analizaba la
serie de armónicos: explorar las leyes materiales del sonido. La
diferencia es que Rameau tenía una pluma; yo tengo un teclado.
Rameau tardaba meses en publicar un tratado; yo tardo milisegundos en
evaluar una función.

El algoritmo Karplus-Strong (1983) es un ejemplo de cómo un modelo
físico puede recrear el sonido de una cuerda frotada usando un delay
con retroalimentación. Kevin Karplus y Alex Strong demostraron que un
ruido blanco filtrado por una línea de retardo con pérdida de energía
produce un sonido sorprendentemente orgánico. No es una grabación: es
una ecuación que vibra.

\begin{verbatim}
// Karplus-Strong con sustain (cuerda frotada - violonchelo)
// La excitación es continua, no un pulso

// Liberar memoria y reiniciar servidor
s.freeAll;
s.quit;
s.boot;

(
SynthDef(\ks_comb, { |out=0, freq=220, amp=0.2, decay=3, pan=0|
    var sig, delayTime;
    delayTime = freq.reciprocal;
    sig = WhiteNoise.ar(amp) * EnvGen.kr(Env.perc(0.001, decay), doneAction: 2);
    sig = CombC.ar(sig, 0.5, delayTime, decay);
    sig = LPF.ar(sig, freq * 4 + 100);
    sig = sig * 2;
    Out.ar(out, Pan2.ar(sig, pan));
}).add;
)

// Probar con una nota
x = Synth(\ks_comb, [\freq, 220, \amp, 0.3, \decay, 4]);
x.free;

// Melodía con sustain
~melodia = Pbind(
    \instrument, \ks_comb,
    \midinote, Pseq([48, 52, 55, 60, 64, 67, 72, 67, 64, 60, 55, 52, 48], inf),
    \dur, Pseq([1, 0.5, 0.5, 1, 0.5, 0.5, 2, 0.5, 0.5, 1, 0.5, 0.5, 2], inf),
    \amp, 0.15,
    \decay, 4
).play;

// Detener
~melodia.stop;
\end{verbatim}

En el laboratorio de Fourier, en cambio, no modelamos una cuerda:
descomponemos el sonido en sus armónicos constituyentes y los
activamos uno a uno. Cada parcial es controlable; el timbre resultante
es la suma de todos ellos. Es la demostración práctica de que el
timbre no es magia: es combinatoria.

\begin{verbatim}
// Laboratorio de Fourier: control de armónicos individuales
// Cada armónico es un parcial que puede activarse/desactivarse
// El timbre resultante es la suma de todos los parciales activos

// Liberar memoria y reiniciar servidor
s.freeAll;
s.quit;
s.boot;

(
SynthDef(\fourier_lab, { |out=0, freq=220, amp=0.08|
    var sig;
    sig = Mix.ar(16.collect { |i|
        var n = i + 1;
        var active = NamedControl.kr(("arm_" ++ n).asSymbol, 1, 0.02);
        SinOsc.ar(freq * n, 0, 1/n) * active
    });
    sig = sig * amp;
    Out.ar(out, sig!2);
}).add;
)

// Probar con fundamental sola (seno puro)
x = Synth(\fourier_lab, [\freq, 220, \amp, 0.15]);

// Activar armónicos uno a uno
x.set(\arm_2, 1);  // 2do armónico (octava)
x.set(\arm_3, 1);  // 3er armónico (quinta)
x.set(\arm_4, 1);  // 4to armónico (doble octava)
x.set(\arm_5, 1);  // 5to armónico (tercera mayor)

// Apagar la fundamental para escuchar solo armónicos
x.set(\arm_1, 0);

// Volver a encenderla
x.set(\arm_1, 1);

// Activar impares para aproximar onda cuadrada
x.set(\arm_7, 1);
x.set(\arm_9, 1);
x.set(\arm_11, 1);

// Liberar
x.free;
\end{verbatim}

Si el laboratorio de Fourier descompone el sonido en armónicos, la
síntesis granular lo fragmenta en partículas temporales: granos de
entre 10 y 50 milisegundos. Cada grano está envuelto en una curva
gaussiana —nace suave, crece hasta un pico, decae suavemente— para
evitar los clics que producirían los cortes bruscos. Miles de granos
superpuestos crean texturas que van desde nubes etéreas hasta paisajes
industriales. Es el equivalente sonoro del puntillismo de Georges
Seurat: en \textbf{Tarde de domingo en la isla de La Grande Jatte}
(1884-1886), los colores no se mezclan en la paleta, sino en el ojo
del espectador. De la suma de puntos diminutos emerge una totalidad
que ninguno de ellos contenía.


\begin{verbatim}
// Síntesis granular: nubes de partículas sonoras
// Cada grano es una onda sinusoidal microscópica
// La textura emerge de la superposición de miles de granos

// Liberar memoria y reiniciar servidor
s.freeAll;
s.quit;
s.boot;

(
SynthDef(\granular, { |out=0, freq=440, amp=0.1, density=50, spread=0.5|
    var sig, trig, grains;
    // Generar granos a una densidad controlada (granos por segundo)
    trig = Dust.ar(density);
    grains = GrainSin.ar(
        numChannels: 2,
        trigger: trig,
        dur: 0.02,                    // Duración de cada grano: 20 ms
        freq: freq * LFNoise1.kr(3).range(1 - spread, 1 + spread),
        pan: LFNoise1.kr(5).range(-1, 1)
    );
    sig = grains * amp * 0.5;
    Out.ar(out, sig);
}).add;
)

// Probar con diferentes densidades
x = Synth(\granular, [\freq, 220, \amp, 0.15, \density, 20, \spread, 0.1]);
x.set(\density, 100);  // Más granos: textura más densa
x.set(\density, 500);  // Muy denso: casi un sonido continuo
x.set(\density, 10);   // Pocos granos: textura punteada
x.set(\spread, 0.5);   // Más variación de frecuencia
x.set(\freq, 440);     // Cambiar frecuencia base

// Liberar
x.free;

// Evolución de una nube granular
(
~nube = Pbind(
    \instrument, \granular,
    \freq, Pseq([110, 220, 330, 440, 330, 220, 110], inf),
    \dur, 4,
    \density, Pseq([10, 50, 200, 500, 200, 50, 10], inf),
    \spread, Pseq([0.05, 0.1, 0.3, 0.5, 0.3, 0.1, 0.05], inf),
    \amp, 0.1
).play;
)

// Detener
~nube.stop;
\end{verbatim}

\begin{verbatim}
library(ggplot2)

set.seed(1884)  # Año de La Grande Jatte

# Crear una cuadrícula de puntos como un lienzo puntillista
n_puntos <- 3000
datos <- data.frame(
  x = runif(n_puntos, 0, 10),
  y = runif(n_puntos, 0, 6),
  # Los colores no están mezclados: son puntos puros
  color = sample(c("#1A5276", "#922B21", "#D4AC0D", "#6C3483", "#2E86C1", "#E74C3C"), 
                 n_puntos, replace = TRUE)
)

ggplot(datos, aes(x = x, y = y, color = color)) +
  geom_point(size = 1.5, alpha = 0.7) +
  scale_color_identity() +
  labs(title = "Puntillismo de Seurat — La Grande Jatte (1884-1886)",
       subtitle = "Cada punto es un color puro. La mezcla ocurre en el ojo, no en el lienzo.") +
  theme_void() +
  theme(plot.background = element_rect(fill = "#FAF0E6", color = NA),
        plot.title = element_text(face = "bold", size = 16, hjust = 0.5),
        plot.subtitle = element_text(color = "gray40", size = 11, hjust = 0.5))
\end{verbatim}

\href{file:///home/spock/Documents/symploke-sonora/img/seurat\_puntillismo.png}{seurat\textsubscript{puntillismo}}

\begin{verbatim}
library(ggplot2)

# Simular una envolvente gaussiana (campana)
t <- seq(-3, 3, length.out = 200)
gaussiana <- exp(-t^2 / 2)

# Simular granos individuales con envolvente gaussiana
granos <- data.frame()
set.seed(42)
for(i in 1:12) {
  centro <- runif(1, -2, 2)
  amplitud <- runif(1, 0.3, 1)
  t_grano <- seq(centro - 0.6, centro + 0.6, length.out = 80)
  grano <- amplitud * exp(-((t_grano - centro)^2) / 0.08)
  granos <- rbind(granos, data.frame(
    tiempo = t_grano,
    amplitud = grano,
    grano = factor(i)
  ))
}

# Gráfico principal: envolvente gaussiana
p1 <- ggplot(data.frame(t = t, y = gaussiana), aes(x = t, y = y)) +
  geom_line(color = "#1A5276", linewidth = 1.5) +
  labs(title = "Envolvente Gaussiana",
       subtitle = "Cada grano respira: nace suave, crece, decae",
       x = "", y = "") +
  theme_minimal(base_size = 12) +
  theme(plot.title = element_text(face = "bold", hjust = 0.5),
        plot.subtitle = element_text(color = "gray40", hjust = 0.5))

# Gráfico de granos
p2 <- ggplot(granos, aes(x = tiempo, y = amplitud, color = grano, group = grano)) +
  geom_line(linewidth = 0.8) +
  scale_color_viridis_d(guide = "none") +
  labs(title = "Granos individuales con envolvente gaussiana",
       subtitle = "12 granos. Superpuestos, crean una textura.",
       x = "Tiempo", y = "Amplitud") +
  theme_minimal(base_size = 12) +
  theme(plot.title = element_text(face = "bold", hjust = 0.5),
        plot.subtitle = element_text(color = "gray40", hjust = 0.5))

# Superponer usando patchwork
library(patchwork)
p1 / p2 +
  plot_annotation(
    title = "Síntesis Granular: la envolvente gaussiana",
    subtitle = "Cada grano es una onda sinusoidal microscópica envuelta en una campana de Gauss",
    theme = theme(plot.title = element_text(face = "bold", size = 16, hjust = 0.5),
                  plot.subtitle = element_text(color = "gray40", size = 11, hjust = 0.5))
  )
\end{verbatim}

\href{file:///home/spock/Documents/symploke-sonora/img/gaussiana\_granular.png}{envolvente gaussiana}
\item Mi praxis: una ecología de herramientas libres
\label{sec:org4c61421}
No trabajo con una herramienta: trabajo con una ecología de
herramientas que se entrelazan. SuperCollider para la síntesis. Csound
para la orquestación. FoxDot para el live coding. Hydrogen para la
percusión. Ardour para la mezcla. LSP y AMS para los efectos.
JackControl para enrutar el audio entre aplicaciones. Emacs para
escribir, documentar, programar, controlar. LilyPond para la notación.
R para los gráficos. Bash para los scripts. El kernel real-time como
base material de todo lo anterior.

Esta ecología no es un capricho: es una necesidad ontológica. Cada
programa es un género discontinuo de materialidad computacional.
Ninguno puede hacerlo todo. Pero entrelazados (\textbf{symploké}), producen
una totalidad que ninguno podría producir por separado.

Mi método es una herencia del jazz: cambio de pantallas, modifico el
código en la marcha, lo llevo al extremo. No parto de la nada:
modifico, forkéo, reapropio. La mayoría de mis códigos son
modificaciones de códigos libres. Como dicta la ética hacker: el
conocimiento no tiene dueño.
\item La música como dato, el dato como música
\label{sec:org599daeb}
En la era digital, la música es un archivo. Un archivo es una
secuencia de bits. Una secuencia de bits es un número. Un número es
cálculo. La música se ha convertido en lo que Pitágoras intuyó y
Fourier demostró: vibración organizada matemáticamente.

Pero hay una diferencia fundamental. Pitágoras creía que el número era
el principio de todas las cosas. El materialismo filosófico afirma lo
contrario: el número es una categoría que emerge de la materia cuando
esta se organiza. El código no es la música: es una representación de
la música. La música es un evento sonoro que se da en el tiempo y en
el espacio; el código es un plano que permite reconstruir ese evento.

Confundir el código con la música es el pecado original del idealismo
digital. Es lo mismo que hizo Pitágoras con el monocordio: confundir
el instrumento con el cosmos. El ordenador no es la música: es una
herramienta para pensar la música. Como el monocordio. Como la
notación de Guido de Arezzo. Como el piano forte.
\item La soberanía tecnológica como elección
\label{sec:orgfcea37a}
En un mundo donde el software privativo controla la producción, la
distribución y el consumo de música, elegir el software libre es un
acto político. No porque el software libre sea "mejor" en términos
técnicos —a veces lo es, a veces no—, sino porque garantiza la
autonomía del usuario. Un músico que usa software privativo depende de
una corporación para producir su propio sonido. Un músico que usa
software libre depende de su propia capacidad de aprender, modificar y
compartir.

Esta elección no es inevitable: es una decisión consciente dentro de
un campo de fuerzas donde el software privativo y el software libre
competen. Yo elijo el software libre. No porque sea la única opción,
sino porque es coherente con una ontología que afirma que el
conocimiento no tiene dueño y que las herramientas que producen
materia sonora deben ser, ellas mismas, libres.
\end{enumerate}
\subsection{Genealogía de la Computación: del ábaco al ordenador}
\label{sec:orga6d2a42}
El ordenador no nació de la nada. Es el resultado de una larga cadena
de invenciones que fueron transformando la relación entre el ser
humano y el cálculo. Cada eslabón de esa cadena es una tecnología que
amplió las posibilidades del pensamiento. Y cada tecnología, como
hemos visto, no es neutral: condiciona lo que se puede pensar y lo que
se puede hacer.
\subsubsection{El ábaco y el algoritmo}
\label{sec:org28cc3b9}
El ábaco —inventado en Mesopotamia hacia el 2700 a.C., difundido por
griegos, romanos, chinos y mayas— es la primera máquina de calcular.
No hace cálculos por sí mismo: externaliza la memoria. Las cuentas en
las varillas representan cantidades que el cerebro puede manipular
visualmente. El ábaco no piensa: hace visible el pensamiento.

En la música, el ábaco tiene un equivalente: el monocordio de
Pitágoras. Ambos son tecnologías de visualización de relaciones
matemáticas. El ábaco muestra números; el monocordio muestra
proporciones. Ambos son precursores del ordenador: máquinas que no
piensan, pero que permiten pensar mejor.
\subsubsection{Leibniz y la máquina de calcular}
\label{sec:org5b89550}
Gottfried Wilhelm Leibniz (1646-1716) no solo desarrolló el cálculo
infinitesimal y definió la música como «aritmética inconsciente del
alma». En 1673 presentó en Londres una máquina de calcular capaz de
sumar, restar, multiplicar y dividir. Era una caja de madera con
ruedas dentadas y una manivela: un cerebro mecánico.

Leibniz soñó con un lenguaje universal —la *characteristica
universalis*— que permitiera expresar todo el conocimiento humano en
símbolos y resolver cualquier disputa mediante el cálculo. Ese sueño
es el precursor directo de la inteligencia artificial. Leibniz no lo
logró, pero su máquina de calcular demostró que el razonamiento podía
ser mecanizado.

\begin{verbatim}
# Máquina de Leibniz (1673): suma, resta, multiplicación, división
# Emulación lógica de la Staffelwalze

maquina_leibniz <- function(a, b, operacion = "suma") {
  switch(operacion,
    "suma" = a + b,
    "resta" = a - b,
    "multiplicacion" = {
      # Leibniz: la multiplicación es suma repetida
      resultado <- 0
      for(i in 1:abs(b)) {
        resultado <- resultado + a
      }
      if(b < 0) resultado <- -resultado
      resultado
    },
    "division" = {
      # Leibniz: la división es resta repetida
      if(b == 0) return(NA)
      cociente <- 0
      resto <- abs(a)
      while(resto >= abs(b)) {
        resto <- resto - abs(b)
        cociente <- cociente + 1
      }
      if((a < 0 && b > 0) || (a > 0 && b < 0)) cociente <- -cociente
      cociente
    }
  )
}

# Demostración
cat("=== Máquina de Leibniz (1673) ===\n")
cat("Suma: 15 + 27 =", maquina_leibniz(15, 27, "suma"), "\n")
cat("Resta: 42 - 19 =", maquina_leibniz(42, 19, "resta"), "\n")
cat("Multiplicación: 6 × 8 =", maquina_leibniz(6, 8, "multiplicacion"), "\n")
cat("División: 144 ÷ 12 =", maquina_leibniz(144, 12, "division"), "\n")
\end{verbatim}

\begin{verbatim}
// Máquina de Leibniz en SuperCollider (1673)
// La Staffelwalze como sintetizador: cada operación produce un sonido
// "La música es aritmética inconsciente del alma" — Leibniz, 1712

s.boot;

// SynthDef: un tono por cada operación
(
SynthDef(\leibniz_tono, { |out=0, freq=440, amp=0.1, duracion=0.3|
    var sig, env;
    env = EnvGen.kr(Env.perc(0.01, duracion, 1, -4), doneAction: 2);
    sig = SinOsc.ar(freq, 0, amp) * env;
    Out.ar(out, sig!2);
}).add;
)

// Ejecutar la máquina de Leibniz
(
~leibniz_sonoro = Routine({
    var a = 15, b = 27;
    var suma = a + b;
    var resta = a - b;
    var multi = a * b;
    var div = if(b != 0) { a / b } { 0 };

    ("=== Máquina de Leibniz (1673) - Sonora ===").postln;
    ("Entrada: a = " ++ a ++ ", b = " ++ b).postln;

    ("Suma: " ++ a ++ " + " ++ b ++ " = " ++ suma).postln;
    Synth(\leibniz_tono, [\freq, (60 + suma).midicps, \amp, 0.1, \duracion, 0.5]);
    0.6.wait;

    ("Resta: " ++ a ++ " - " ++ b ++ " = " ++ resta).postln;
    Synth(\leibniz_tono, [\freq, (60 + resta.abs).midicps, \amp, 0.08, \duracion, 0.5]);
    0.6.wait;

    ("Multiplicación: " ++ a ++ " × " ++ b ++ " = " ++ multi).postln;
    Synth(\leibniz_tono, [\freq, (60 + (multi % 40)).midicps, \amp, 0.06, \duracion, 0.5]);
    0.6.wait;

    ("División: " ++ a ++ " ÷ " ++ b ++ " = " ++ div).postln;
    Synth(\leibniz_tono, [\freq, (60 + (div * 40).asInteger).midicps, \amp, 0.05, \duracion, 0.5]);
    0.6.wait;

    ("Leibniz: 'La música es aritmética inconsciente del alma'").postln;
});

~leibniz_sonoro.play;
)
\end{verbatim}
\subsubsection{Babbage y la máquina analítica}
\label{sec:orgc2b4404}
Charles Babbage (1791-1871), matemático británico, diseñó dos máquinas
que nunca llegó a construir completamente: la máquina diferencial
(1822), para calcular tablas matemáticas sin errores humanos, y la
máquina analítica (1837), capaz de ejecutar cualquier tipo de cálculo
mediante instrucciones programadas en tarjetas perforadas.

La máquina analítica era, en esencia, un ordenador mecánico. Tenía una
unidad de procesamiento (el «molino»), una memoria (el «almacén») y un
sistema de entrada/salida basado en las tarjetas perforadas del telar
de Jacquard. Babbage no la construyó por falta de fondos y por
conflictos con su ingeniero jefe, pero su diseño era correcto. En
1991, el Museo de Ciencias de Londres construyó una réplica funcional:
la máquina funciona.
\subsubsection{Ada Lovelace: la primera programadora}
\label{sec:org2bbc863}
Ada Lovelace (1815-1852), hija del poeta Lord Byron, fue la primera
persona en comprender que la máquina analítica de Babbage no era solo
una calculadora: era una máquina de propósito general capaz de
manipular símbolos, no solo números. En sus notas a la traducción de
un artículo sobre la máquina (1843), Lovelace escribió el primer
algoritmo diseñado para ser ejecutado por una máquina: un método para
calcular los números de Bernoulli.

Lovelace intuyó algo que tardaríamos un siglo en comprender
plenamente: que una máquina de calcular podía convertirse en una
máquina de crear. Escribió: «La máquina analítica no tiene
pretensión alguna de originar nada. Puede hacer cualquier cosa que
sepamos ordenarle que haga». Es la definición más precisa de lo que
hoy llamamos un programa.
\subsubsection{Turing y la máquina universal}
\label{sec:orgc0556f4}
Alan Turing (1912-1954), matemático británico, formalizó en 1936 el
concepto de «máquina universal»: un dispositivo teórico capaz de
simular cualquier otra máquina. La máquina de Turing no era un objeto
físico, sino un modelo matemático: una cinta infinita, un cabezal que
lee y escribe símbolos, un conjunto de reglas. Esa abstracción es el
fundamento de todos los ordenadores modernos.

Durante la Segunda Guerra Mundial, Turing trabajó en Bletchley Park
descifrando los códigos de la máquina Enigma, utilizada por los nazis
para encriptar sus comunicaciones. La máquina Enigma —un dispositivo
electromecánico con rotores y conexiones— era una tecnología de
ocultación. Turing construyó una tecnología de revelación: la
«Bombe», un ordenador primitivo que automatizaba el descifrado.

La relación entre Turing y la música es indirecta pero profunda.
Turing no era músico, pero su concepto de «máquina universal» es el
fundamento de todos los sintetizadores digitales. SuperCollider, Pure
Data, Csound: todos son máquinas de Turing aplicadas al sonido.
\subsubsection{Von Neumann y la arquitectura del ordenador moderno}
\label{sec:org5b6c1eb}
John von Neumann (1903-1957), matemático húngaro-estadounidense,
definió en 1945 la arquitectura que siguen usando todos los
ordenadores modernos: una unidad de procesamiento, una memoria que
almacena tanto datos como instrucciones, y un bus que las conecta. La
arquitectura von Neumann permitió que los programas no fueran cables,
sino datos: se podían modificar, copiar, distribuir.

Esa flexibilidad es la que permite que un mismo ordenador pueda
ejecutar un procesador de texto, un navegador de internet o un
sintetizador. La arquitectura von Neumann es, en el fondo, una
ontología: la materia computacional (\(M_1\)) es indiferente a lo que
representa. Los mismos bits que codifican una sinfonía de Beethoven
pueden codificar una imagen o un texto. El significado no está en la
materia: está en la interpretación (\(M_3\)).
\subsubsection{Del ordenador al sintetizador}
\label{sec:orgd7b4dca}
En 1957, Max Mathews escribió MUSIC I en un IBM 704: el primer
programa de síntesis digital. El ordenador tardaba horas en generar
unos segundos de sonido. Pero el principio estaba establecido: el
cálculo podía convertirse en vibración.

En los años ochenta, Miller Puckette creó Pure Data, un entorno de
programación gráfica para síntesis de audio. En los años noventa,
James McCartney creó SuperCollider, un lenguaje de programación
diseñado específicamente para la síntesis y el procesamiento de audio
en tiempo real.

El ordenador ya no es una herramienta: es un instrumento. Y el músico
que programa no es un técnico: es un filósofo que explora las leyes
materiales del sonido mediante código.
\subsubsection{La ética hacker y la música libre}
\label{sec:org11d2377}
En 1971, Richard Stallman —entonces un joven programador del MIT— se
topó con una impresora Xerox que se atascaba constantemente. Stallman
quiso modificar el código para que la impresora avisara a los usuarios
cuando se atascaba. Xerox se negó a darle el código fuente.

Ese gesto —negar el acceso al código— es el pecado original de la era
digital. Stallman respondió fundando la Free Software Foundation y
escribiendo la GNU General Public License: un documento legal que
garantiza cuatro libertades fundamentales: usar, estudiar, redistribuir
y mejorar el software.

La ética hacker —acceso libre, desconfianza de la autoridad, mérito
basado en la calidad del trabajo— es la aplicación de esa ética al
mundo digital. Y en la música, el software libre no es una preferencia
técnica: es una posición ontológica. Si la música es materia sonora,
las herramientas que la producen deben ser libres. No por ideología,
sino por coherencia materialista: el conocimiento no tiene dueño.
\subsubsection{Brian Ferneyhough y la Nueva Complejidad: un espejo incómodo}
\label{sec:orge0371f7}
Brian Ferneyhough (1943-) es la figura más visible de la Nueva
Complejidad. Británico, formado en el Royal College of Music, profesor
en Friburgo, La Haya, Stanford y San Diego. Ha recibido el Premio
Ernst von Siemens (2007), dotado con 200,000 euros, uno de los
galardones más prestigiosos de la música contemporánea europea. Su
música se programa regularmente en Donaueschingen, Darmstadt e IRCAM,
instituciones financiadas con fondos públicos.
\begin{enumerate}
\item La partitura como campo de fuerzas
\label{sec:org6b2784b}
Ferneyhough escribe partituras que ningún intérprete puede ejecutar
con precisión absoluta. Ritmos anidados —tresillos dentro de
quintillos dentro de septillos—, microtonos, dinámicas que cambian
nota a nota, articulaciones obsesivamente especificadas. La partitura
no es una instrucción: es un campo de fuerzas. El intérprete debe
negociar con la imposibilidad.

Desde el materialismo filosófico, Ferneyhough representa el idealismo
musical llevado a su paroxismo. La partitura (\(M_3\)) se independiza
del sonido (\(M_1\)) y de la percepción (\(M_2\)). El compositor no espera
que su música sea ejecutada tal como la escribió: espera que la
tensión entre lo escrito y lo ejecutable genere una experiencia
estética. La idea desborda la materia.

\begin{verbatim}
\version "2.24.0"

\score {
  \new Staff {
    \time 4/4
    \override Score.BarNumber.break-visibility = ##(#f #f #f)
    \cadenzaOn
    % Tresillo dentro de quintillo: la marca de la Nueva Complejidad
    \tuplet 5/4 {
      c''8 \tuplet 3/2 { d''16 e'' f'' } g'' a''
    }
    \tuplet 7/8 {
      b''16 c''' d''' e''' f''' g''' a'''
    }
    \tuplet 5/4 {
      \tuplet 3/2 { c''16 d'' e'' } \tuplet 3/2 { f'' g'' a'' } b'' c'''
    }
    \bar "|."
  }
  \layout {
    indent = 0
    ragged-right = ##f
  }
}
\end{verbatim}

\begin{figure}[htbp]
\centering
\includegraphics[width=.9\linewidth]{/home/spock/Documents/symploke-sonora/img/nueva_complejidad_ritmos.png}
\caption{Ritmos anidados al estilo de la Nueva Complejidad}
\end{figure}

\begin{verbatim}
// Nueva Complejidad: polirritmia extrema en SuperCollider
// Cinco voces independientes, cada una con su propia lógica temporal
// Inspirado en Ferneyhough: la partitura como campo de fuerzas

s.boot;

(
// SynthDef simple para cada voz
SynthDef(\voz_compleja, { |out=0, freq=440, amp=0.08, gate=1, pan=0|
    var sig, env;
    env = EnvGen.kr(Env.perc(0.01, 0.3, 1, -4), gate, doneAction: 2);
    sig = SinOsc.ar(freq, 0, amp) * env;
    Out.ar(out, Pan2.ar(sig, pan));
}).add;
)

(
// Cinco voces con patrones rítmicos independientes
~complejidad = Routine({
    var voces = 5;
    var patrones = [
        // Voz 1: quintillos (5 notas en espacio de 4)
        Pbind(\instrument, \voz_compleja,
            \midinote, Pseq([72, 74, 76, 77, 79], inf),
            \dur, 0.16, \amp, 0.06, \pan, -0.8),
        // Voz 2: septillos (7 notas en espacio de 4)
        Pbind(\instrument, \voz_compleja,
            \midinote, Pseq([60, 62, 64, 65, 67, 69, 71], inf),
            \dur, 0.114, \amp, 0.05, \pan, -0.4),
        // Voz 3: tresillos (3 notas en espacio de 2)
        Pbind(\instrument, \voz_compleja,
            \midinote, Pseq([48, 52, 55, 60, 64, 67], inf),
            \dur, 0.333, \amp, 0.06, \pan, 0),
        // Voz 4: onceillos (11 notas en espacio de 8)
        Pbind(\instrument, \voz_compleja,
            \midinote, Pseq([67, 65, 64, 62, 60, 59, 57, 55, 53, 52, 50], inf),
            \dur, 0.145, \amp, 0.04, \pan, 0.4),
        // Voz 5: treceavos (13 notas en espacio de 8)
        Pbind(\instrument, \voz_compleja,
            \midinote, Pseq([36, 38, 40, 41, 43, 45, 47, 48, 50, 52, 53, 55, 57], inf),
            \dur, 0.123, \amp, 0.05, \pan, 0.8)
    ];

    ("=== Nueva Complejidad: polirritmia extrema ===").postln;
    ("Cinco voces, cinco lógicas temporales independientes").postln;
    ("Ningún intérprete humano puede ejecutar esto con precisión").postln;

    // Iniciar todas las voces
    patrones.do { |patron, i|
        ("Voz " ++ (i+1) ++ " iniciada").postln;
        patron.play;
    };

    30.wait;  // Dejar sonar durante 30 segundos

    ("La partitura como campo de fuerzas: imposible, pero real.").postln;
});

~complejidad.play;
)

// Para detener: Cmd + .
\end{verbatim}

\begin{verbatim}
// Nueva Complejidad: polirritmia extrema en SuperCollider
// Cinco voces independientes, cada una con su propia lógica temporal
// Inspirado en Ferneyhough: la partitura como campo de fuerzas

s.boot;

(
// SynthDef simple para cada voz
SynthDef(\voz_compleja, { |out=0, freq=440, amp=0.08, gate=1, pan=0|
    var sig, env;
    env = EnvGen.kr(Env.perc(0.01, 0.3, 1, -4), gate, doneAction: 2);
    sig = SinOsc.ar(freq, 0, amp) * env;
    Out.ar(out, Pan2.ar(sig, pan));
}).add;
)

(
// Cinco voces con patrones rítmicos independientes
~complejidad = Routine({
    var voces = 5;
    var patrones = [
        // Voz 1: quintillos (5 notas en espacio de 4)
        Pbind(\instrument, \voz_compleja,
            \midinote, Pseq([72, 74, 76, 77, 79], inf),
            \dur, 0.16, \amp, 0.06, \pan, -0.8),
        // Voz 2: septillos (7 notas en espacio de 4)
        Pbind(\instrument, \voz_compleja,
            \midinote, Pseq([60, 62, 64, 65, 67, 69, 71], inf),
            \dur, 0.114, \amp, 0.05, \pan, -0.4),
        // Voz 3: tresillos (3 notas en espacio de 2)
        Pbind(\instrument, \voz_compleja,
            \midinote, Pseq([48, 52, 55, 60, 64, 67], inf),
            \dur, 0.333, \amp, 0.06, \pan, 0),
        // Voz 4: onceillos (11 notas en espacio de 8)
        Pbind(\instrument, \voz_compleja,
            \midinote, Pseq([67, 65, 64, 62, 60, 59, 57, 55, 53, 52, 50], inf),
            \dur, 0.145, \amp, 0.04, \pan, 0.4),
        // Voz 5: treceavos (13 notas en espacio de 8)
        Pbind(\instrument, \voz_compleja,
            \midinote, Pseq([36, 38, 40, 41, 43, 45, 47, 48, 50, 52, 53, 55, 57], inf),
            \dur, 0.123, \amp, 0.05, \pan, 0.8)
    ];

    ("=== Nueva Complejidad: polirritmia extrema ===").postln;
    ("Cinco voces, cinco lógicas temporales independientes").postln;
    ("Ningún intérprete humano puede ejecutar esto con precisión").postln;

    // Iniciar todas las voces
    patrones.do { |patron, i|
        ("Voz " ++ (i+1) ++ " iniciada").postln;
        patron.play;
    };

    30.wait;  // Dejar sonar durante 30 segundos

    ("La partitura como campo de fuerzas: imposible, pero real.").postln;
});

~complejidad.play;
)

// Para detener: C- c C-s
\end{verbatim}
\item Un espejo incómodo
\label{sec:org149f559}
Ferneyhough es, para este ensayo, un espejo incómodo. Representa todo
lo que la academia europea valora: complejidad, especialización,
innovación formal, financiación institucional. Y representa todo lo
que a este ensayo le falta: títulos, reconocimiento, acceso a los
circuitos de legitimación.

Pero el espejo también refleja lo contrario. Ferneyhough escribe
música que no se puede tocar. Este ensayo describe música que se toca
sin necesidad de ser escrita. Ferneyhough trabaja dentro de la
academia. Este ensayo trabaja desde fuera, intentando entrar.
Ferneyhough confía en \(M_3\) (el concepto) y desconfía de \(M_2\) (la
percepción). Este ensayo confía en \(M_1\) (la materia sonora) y en
\(M_2\) (la experiencia incorporada), y considera que \(M_3\) debe estar
al servicio de ambas.

La pregunta no es quién tiene razón. La pregunta es por qué la
complejidad de Ferneyhough es financiada con fondos públicos europeos
y la complejidad de un músico que ha atravesado retenes paramilitares
para enseñar a leer partituras en pueblos sin presencia estatal no
recibe apoyo alguno. La respuesta no es musical: es política. La
academia financia lo que se parece a ella.
\item Contrafactual
\label{sec:org6c42458}
¿Qué tendría que ocurrir para que la academia colombiana reconociera
como Sabedor a alguien que no pasó por sus aulas? Que la academia
dejara de definir el saber por los títulos y empezara a definirlo por
la práctica. Que la universidad aplicara la ley que ya existe —la Ley
30 de 1992, la figura del Sabedor— en lugar de ignorarla. Que la
complejidad de la vida —la que se aprende en las zonas de conflicto,
en los barrios marginados, en las comunidades de software libre—
fuera considerada tan valiosa como la complejidad de las partituras.

Eso es lo que este ensayo solicita. No un privilegio: el cumplimiento
de una ley que ya está vigente. No una excepción: la aplicación de un
principio que la propia academia latinoamericana ha validado a través
de la Investigación-Acción Participativa de Fals Borda. No un favor:
justicia epistémica.
\end{enumerate}
\section{De la ontología a la educación: ¿para qué saber?}
\label{sec:org22941d8}
Este ensayo ha recorrido un largo camino. Desde los presocráticos hasta
Ferneyhough. Desde el monocordio de Pitágoras hasta SuperCollider.
Desde las zonas de conflicto del Ariari hasta el software libre. Desde
la IAP de Fals Borda hasta la ética hacker. Ha demostrado —o ha
intentado demostrar— que la música no es sentimiento, sino materia
organizada. Que la discontinuidad de la materia no es caos, sino
condición de posibilidad del orden. Que la \textbf{symploké} no es un
concepto abstracto, sino una biografía.

Pero todo este recorrido tiene un propósito que trasciende la
ontología. Porque la pregunta no es solo "¿qué es la música?". La
pregunta es también "¿para qué saber?".
\subsection{El diagnóstico: una educación fragmentada}
\label{sec:org262c2e0}

Vivimos una paradoja. Nunca hubo tanto acceso a la información. Nunca
hubo tan poca capacidad de entrelazarla. La educación primaria y
secundaria sigue organizada en asignaturas estancas: matemáticas por
un lado, música por otro, filosofía por otro. La universidad
profundiza la fragmentación: el músico no estudia física, el físico no
estudia música, el filósofo no programa. El resultado es una doble
pobreza: la ciencia recurre a metáforas místicas porque ha perdido el
oído, y el arte se refugia en el mercado o en la academia porque ha
perdido el mundo.

Estanislao Zuleta diagnosticó este problema en su \textbf{Elogio de la
dificultad}: la educación no enseña a pensar, enseña a obedecer. No
forma ciudadanos críticos, sino consumidores dóciles. Leer no es
recibir información: es trabajar. Escribir no es transcribir: es
organizar el pensamiento. Y sin embargo, la escuela trata la lectura
como un ejercicio mecánico y la escritura como una habilidad técnica.
Se lee sin comprender. Se escribe sin pensar.

Jesús Maestro, el cervantista español, ha llevado esta crítica más
lejos: se prueba mejor la inteligencia interpretando un soneto de
Quevedo que resolviendo ejercicios lógico-matemáticos. No porque las
matemáticas sean inútiles —este ensayo ha demostrado lo contrario—,
sino porque el soneto exige sensibilidad, contexto, historia,
ambigüedad. La educación actual entrena para el examen. No para la
vida.
\subsection{La propuesta: el \textbf{trivium} y el \textbf{quadrivium} como modelo}
\label{sec:org9212da3}

Este ensayo no propone una reforma educativa. Pero sí propone un
modelo que ha demostrado su eficacia durante siglos: las siete artes
liberales. El \textbf{trivium} —gramática, lógica, retórica— enseña a pensar.
El \textbf{quadrivium} —aritmética, geometría, música, astronomía— enseña a
conocer. Juntas, forman una educación integral que no separa las
ciencias de las artes, ni la teoría de la práctica, ni el conocimiento
de la experiencia.

En la era del silicio, estas siete artes no han perdido vigencia: la
han recuperado. Programar es ejercer la gramática y la lógica sobre un
lenguaje formal. Diseñar un sintetizador es aplicar la aritmética y la
geometría a las ondas sonoras. Componer es hacer música, que sigue
siendo número en el tiempo. Y explorar el universo digital es una
forma de astronomía: no miramos las estrellas, sino los datos; pero la
estructura es la misma.
\subsection{El método: aprender haciendo, aprender transformando}
\label{sec:org3457675}

La Investigación-Acción Participativa de Orlando Fals Borda ha sido el
blindaje metodológico de este ensayo. Pero la IAP no es solo una
metodología de investigación: es una filosofía de la educación. No se
aprende desde fuera. Se aprende haciendo, transformando, siendo
transformado. Como yo aprendí con los músicos naturales del Ariari. Mi
tarea no era enseñarles: era encontrar estrategias para que ellos
tuvieran herramientas sin perder su esencia. Ellos me enseñaron tanto
como yo pude ofrecerles.

Esta reciprocidad es lo que falta en la educación formal. El profesor
enseña, el alumno aprende. Pero ¿cuándo aprende el profesor? ¿Cuándo
la academia reconoce que el saber puede venir de fuera? La figura del
Sabedor —contemplada en la legislación colombiana pero raramente
aplicada— es un intento de responder a esa pregunta. No todos los que
saben pasaron por la universidad. La ley lo reconoce. La práctica lo
ignora.
\subsection{La infraestructura: el software libre como condición material}
\label{sec:org7777990}

Si el conocimiento es libre, las herramientas para producirlo también
deben serlo. No por ideología, sino por coherencia materialista. Un
músico que usa software privativo depende de una corporación para
producir su propio sonido. Un estudiante que solo conoce Windows no
sabe cómo funciona su ordenador. Un profesor que usa plataformas
comerciales para dar clase está cediendo el control de su aula a una
empresa.

El software libre —Emacs, SuperCollider, LilyPond, R, Bash— no es una
preferencia técnica: es una infraestructura de soberanía. Permite
aprender sin pedir permiso. Crear sin pagar licencias. Compartir sin
violar leyes. En un mundo donde el conocimiento está cada vez más
controlado por corporaciones, aprender a programar con herramientas
libres es un acto de resistencia.
\subsection{¿Para qué saber?}
\label{sec:orgabbd094}

Este ensayo no tiene una respuesta definitiva. Pero tiene una apuesta.

Saber para entrelazar. Para que la música dialogue con las
matemáticas, la filosofía con la computación, la academia con la
calle. Para que un niño de Aguablanca pueda leer una partitura como
lee un libro. Para que la complejidad de la vida —la que se aprende en
las zonas de conflicto, en los barrios marginados, en las comunidades
de software libre— sea considerada tan valiosa como la complejidad de
las aulas.

Saber para no pedir permiso, pero tampoco para imponer. Porque la
academia no es dueña del saber: es su custodia. Y aunque a veces
olvide su misión, sigue siendo el lugar donde el conocimiento puede
ser discutido, confrontado, enriquecido. Mi crítica no es un ataque:
es una invitación. Mi reclamo no es una ruptura: es una propuesta de
diálogo. Vengo de fuera, sí. Pero vengo con respeto. Con argumentos.
Con código. Con música. Con la esperanza de que la academia —mi
interlocutora, no mi enemiga— sepa escuchar.

Saber, en fin, para sembrar con otros. Como dice la sabiduría del
pueblo nasa del Cauca colombiano: «Caminar la palabra es sembrar el
pensamiento. Quien camina solo, siembra solo. Quien camina con otros,
siembra con otros.» Este ensayo ha sido una larga caminata. No la he
hecho solo. La he hecho con Panchis, con Macky, con la Gaitana, con
los músicos naturales del Ariari, con los compañeros de Asilo 38, con
los Misioneros Claretianos, con los estudiantes que me preguntaban
para qué aprender a leer partituras. Y también con una máquina de
predecir palabras que me ayudó a ordenar el pensamiento. A todos ellos
les debo lo que aquí está escrito. Las imperfecciones son mías. Los
aciertos, de todos.
\section{Petición: Reconocimiento del Saber Situado}
\label{sec:orga62d8e7}
Este ensayo ha sido, desde el principio, una petición.

No pido un título formal. No aspiro a que la academia me conceda lo
que no tengo. Aspiro a que reconozca lo que ya existe: un saber
producido fuera de sus aulas, pero con el mismo rigor, la misma
densidad y la misma pertinencia que el saber que se produce dentro de
ellas.

La figura del Sabedor —contemplada en los reglamentos internos de la
Universidad de Caldas, en consonancia con la figura del Experto
definida por la Ley 30 de 1992 y la Sentencia C-109 de 1994— es la
vía para ese reconocimiento. No es un privilegio: es un derecho. No es
una excepción: es una figura reglamentaria que raramente se aplica. Lo
que solicito no es que se invente nada nuevo, sino que se cumpla lo
que ya está escrito.

Mi trayectoria acredita:

\begin{itemize}
\item Preparación especial en el campo de la filosofía de la música,
la ontología materialista y la praxis con software libre, como
demuestran las páginas precedentes.
\item Preparación empírica: experiencia en territorios sin presencia
efectiva del Estado, trabajo con músicos naturales, producción
musical en contextos de vulnerabilidad, participación en el
movimiento de Resistencia Cultural.
\item Preparación técnica: dominio de instrumentos, arreglos, producción
discográfica, y una ecología de herramientas de software libre
(SuperCollider, Csound, LilyPond, Emacs/Lisp, R, Bash) documentada
en los códigos que acompañan este ensayo.
\end{itemize}

Solicito, por tanto, que este documento sea considerado como base para
el reconocimiento de mis saberes bajo la figura del Sabedor que
contemplan los reglamentos de la Universidad de Caldas, y que un
comité curricular evalúe mi trayectoria para convalidar mis
aprendizajes previos y trazar una ruta formativa continua dentro del
Departamento de Filosofía.

No pido una respuesta inmediata. Pido una lectura. El resto —el
debate, la crítica, la decisión— corresponde a la academia. Yo ya he
hecho mi parte.
\section{Conclusión}
\label{sec:org9a52c6f}
Este ensayo comenzó con una pregunta: ¿cómo se construye conocimiento
musical cuando el camino académico se interrumpe, no por fracaso, sino
por la fuerza de la propia experiencia? La respuesta, desplegada a lo
largo de estas páginas, es: entrelazando. Entrelazando la academia con
la calle, la filosofía con la computación, la música con las
matemáticas, la tradición escrita con la tradición oral, el carbono
con el silicio.

No he pretendido escribir un tratado sistemático. He pretendido
mostrar que el materialismo filosófico —la \textbf{symploké}, la
discontinuidad de la materia, el no-monismo— no es una abstracción,
sino una forma de vida. Una forma de organizar la experiencia. Una
forma de hacer música.

Queda mucho por hacer. Las lecciones de música y algoritmos que este
ensayo anuncia tendrán que esperar a un futuro próximo. Los códigos
que aquí se presentan no son un punto de llegada, sino un punto de
partida. La conversación con el Grupo Tántalo —si llega a ocurrir— no
será un examen, sino un diálogo entre pares.

Como dijo Estanislao Zuleta: la dificultad no es un obstáculo, sino el
método mismo del conocimiento auténtico. Y como dice la sabiduría del
pueblo nasa: caminar la palabra es sembrar el pensamiento. Este ensayo
ha sido una larga caminata. No la he hecho solo. Las imperfecciones
son mías. Los aciertos, de todos.
\section{Agradecimientos}
\label{sec:orgb380365}
A mi madre, Beatriz, y a mi padre, Guillermo: me dieron la vida y el
nombre. A mis hijos —Amada, Felicidad y Nicolás—: me dieron el futuro.
A Gloria, madre de mis hijos: me dio la familia que somos. A mi
hermana Antares y a mi hermano Iván, mi hermano mayor: me dieron la
infancia. A mi tía Geovanny y a su esposo Polo: me dieron el ánimo y
los contactos. A mi tío Salomón: cuando perdí todos mis instrumentos,
él me dio los recursos para volver a tenerlos. Sin esos instrumentos,
no habría música. Sin música, no habría libro.

A mis maestros de música: Andrés Sánchez, Carlos Julio Ramírez, Jaime
Cardona, Mario Posada, Edy Martínez. A todos los licenciados y
maestros que me formaron en la Orquesta Sinfónica Juvenil de Bogotá,
en la Universidad INCA, en la Universidad de Cundinamarca.

A los músicos que caminaron conmigo: Panchis, Macky, MC La Gaitana,
César Reyes. A la Tropa Malagamba. A Ganjariquies. A la Resistencia
Cultural. A Asilo 38. A Skartel. A los músicos naturales que me
enseñaron que la música no se impone: se entrelaza.

A los Misioneros Claretianos, salvoconducto para entrar a territorios
donde el Estado no hacía presencia.

A mis maestros de pensamiento: Estanislao Zuleta, Gustavo Bueno,
Vicente Chuliá, Orlando Fals Borda, Thomas Kuhn, Orlando Mejía Rivera.

A la comunidad del software libre: Richard Stallman, los
desarrolladores de GNU Emacs, SuperCollider, LilyPond, R, Csound,
Debian GNU/Linux. Sin sus herramientas, este libro no existiría.

A DeepSeek, modelo de lenguaje con el que dialogué durante semanas.
No piensa, no siente, no comprende. Pero procesa. Y en ese
procesamiento, me obligó a formular mis ideas con claridad.

A todas las personas que me preguntaron: "¿para qué sirve leer y
escribir música?". Este libro es mi respuesta.

A la Universidad de Caldas, al Departamento de Filosofía, al Grupo de
Investigación Tántalo. Si este libro llega a sus manos, espero que
sepan ver en él no un reclamo, sino una invitación.

El proceso completo de escritura, los códigos fuente, las versiones
intermedias y los archivos de este documento pueden consultarse en el
repositorio público:

\url{https://github.com/rigelsimploke/symploke-sonora}
\section{Referencias}
\label{sec:orgecdf1de}
\subsection{Legislación y jurisprudencia}
\label{sec:orgc970112}
\begin{itemize}
\item {[}@ley30-1992] Ley 30 de 1992. \textbf{Por la cual se organiza el servicio
público de la Educación Superior}. Congreso de la República de
Colombia.
\item {[}@sentencia-c109-1994] Corte Constitucional de Colombia. \textbf{Sentencia
C-109 de 1994}.
\end{itemize}
\subsection{Filosofía y materialismo}
\label{sec:org3024c8f}
\begin{itemize}
\item {[}@aristoteles-categorias] Aristóteles. \textbf{Categorías} (Κατηγορίαι).
\item {[}@bueno1992] Bueno, Gustavo. \textbf{Teoría del cierre categorial}.
Pentalfa, 1992.
\item {[}@chulia-musica] Chuliá, Vicente. \textbf{Materialismo filosófico y
música}.
\item {[}@fgbueno-pluralismo] Fundación Gustavo Bueno. \textbf{Los tres pilares del
pluralismo de Gustavo Bueno}.
\item {[}@fgbueno-particula-dios] Fundación Gustavo Bueno. \textbf{Crítica a la
"partícula de Dios"}.
\end{itemize}
\subsection{Estética y teoría del arte}
\label{sec:org7a2e24c}
\begin{itemize}
\item {[}@baumgarten1750] Baumgarten, Alexander Gottlieb. \textbf{Aesthetica}.
\begin{enumerate}
\item 
\end{enumerate}
\item {[}@kant1790] Kant, Immanuel. \textbf{Crítica del Juicio}. 1790.
\end{itemize}
\subsection{Epistemología e historia de la ciencia}
\label{sec:org143ada3}
\begin{itemize}
\item {[}@kuhn1962] Kuhn, Thomas. \textbf{La estructura de las revoluciones
científicas}. 1962.
\item {[}@chomsky2023] Chomsky, Noam; Roberts, Ian; Watumull, Jeffrey. \textbf{The
False Promise of ChatGPT}. The New York Times, 8 de marzo de 2023.
\end{itemize}
\subsection{Música y organología}
\label{sec:orge07ea0d}
\begin{itemize}
\item {[}@hornbostel-sachs1914] Hornbostel, Erich M. von; Sachs, Curt.
\textbf{Systematik der Musikinstrumente}. 1914.
\item {[}@conard2009] Conard, N.J., Malina, M., \& Münzel, S.C. (2009). New
flutes document the earliest musical tradition in southwestern
Germany. \textbf{Nature}, 460, 737-740.
\item {[}@turk1997] Turk, I. (1997). \textbf{Mousterian Bone Flute?}. Ljubljana:
Založba ZRC.
\end{itemize}
\subsection{Pensamiento latinoamericano}
\label{sec:org976ebe9}
\begin{itemize}
\item {[}@zuleta1980] Zuleta, Estanislao. \textbf{Elogio de la dificultad}. 1980.
\item {[}@falsborda] Fals Borda, Orlando. \textbf{Ciencia propia y colonialismo
intelectual}. 1970.
\item {[}@mejia-medicina] Mejía Rivera, Orlando. \textbf{La medicina arcaica}.
\end{itemize}
\subsection{Software libre}
\label{sec:org023038f}
\begin{itemize}
\item {[}@stallman2002] Stallman, Richard. \textbf{Free Software, Free Society}.
\begin{enumerate}
\item 
\end{enumerate}
\end{itemize}
\section{Anexos}
\label{sec:orge5c6591}
\begin{itemize}
\item[{$\square$}] Copia de la Ley 30 de 1992 (artículos pertinentes).
\item[{$\square$}] Copia de la Sentencia C-109 de 1994.
\item[{$\square$}] Códigos fuente (disponibles en el repositorio GitHub).
\end{itemize}
\end{document}
